% Definitive paper source for Erd\H{o}s Problem 536.
\documentclass[11pt]{article}

\usepackage[a4paper,margin=1in]{geometry}
\usepackage[T1]{fontenc}
\usepackage[utf8]{inputenc}
\usepackage{lmodern}
\usepackage{microtype}
\usepackage{amsmath,amssymb,amsthm,mathtools}
\usepackage{booktabs}
\usepackage{xurl}
\usepackage{authblk}
\usepackage{enumitem}
\usepackage{hyperref}

\allowdisplaybreaks
\numberwithin{equation}{section}

\hypersetup{
  unicode=true,
  colorlinks=true,
  linkcolor=blue,
  citecolor=blue,
  urlcolor=blue,
  hypertexnames=false,
  pdftitle={A Proposed Complete Solution to Erdős Problem 536},
  pdfauthor={Shouqiao Wang}
}
\urlstyle{same}

\theoremstyle{plain}
\newtheorem{theorem}{Theorem}[section]
\newtheorem{proposition}[theorem]{Proposition}
\newtheorem{lemma}[theorem]{Lemma}
\newtheorem{corollary}[theorem]{Corollary}

\newcommand{\F}{\mathbb F}
\newcommand{\Pp}{\mathbb P}
\newcommand{\E}{\mathbb E}
\newcommand{\1}{\mathbf 1}
\newcommand{\dd}{\,\mathrm d}
\newcommand{\e}{\mathrm e}
\newcommand{\eps}{\varepsilon}
\newcommand{\cB}{\mathcal B}
\newcommand{\cM}{\mathcal M}
\newcommand{\cR}{\mathcal R}
\newcommand{\supp}{\operatorname{supp}}
\newcommand{\lcm}{\operatorname{lcm}}
\newcommand{\TV}{\operatorname{TV}}
\newcommand{\kappacap}{\kappa_{\mathrm{cap}}}

\title{
A Proposed Complete Solution to Erd\H{o}s Problem 536
}
\author{Shouqiao Wang}
\affil{Columbia University\hspace{2.5em}Multiscalar Intelligence}
\date{}

\begin{document}
\maketitle

\begin{abstract}
Let $f(N)$ be the largest size of a set $A\subseteq\{1,\ldots,N\}$
containing no three distinct integers whose three pairwise least common
multiples are equal.  We prove that $f(N)=o(N)$, answering a question of
Erd\H{o}s.  The proof first reduces positive density to a weighted
finite-prime extremal problem and then deletes all prime exponents equal to
one.  The resulting squarefree problem is controlled by balanced
pair-product cubes and the cap-set theorem.  A five-state prime-band model,
an exact factorial insertion identity, and matching first- and second-moment
estimates produce cube laws whose word marginals are asymptotically flat.
The only external mathematical inputs are standard prime-number estimates,
the Brun--Titchmarsh inequality, and the polynomial-method cap-set bound.
Finite algebraic and combinatorial checks are reproduced by a companion
Python verifier.
\end{abstract}

\section{Introduction}\label{sec:introduction}

Call a set of positive integers \emph{safe} if it contains no three
distinct members $a,b,c$ satisfying
\begin{equation}\label{eq:forbidden}
 \lcm(a,b)=\lcm(a,c)=\lcm(b,c).
\end{equation}
Let $f(N)$ be the maximum size of a safe subset of $[N]=\{1,\ldots,N\}$.
Erd\H{o}s asked whether every set of positive upper density must contain a
triple of the form \eqref{eq:forbidden} \cite[p.~646]{Erdos1964}; see also
the modern formulation \cite{Bloom536}.  Abbott and Gardner gave an early
lower-bound construction for the associated extremal problem
\cite{AbbottGardner}.

Our main result is the following.

\begin{theorem}\label{thm:main}
As $N\to\infty$,
\[
 f(N)=o(N).
\]
\end{theorem}

Here is the proof strategy.  Equal pairwise least common multiples have an
exact pair-product form
\[
 txy,\qquad txz,\qquad tyz,
\]
with $x,y,z$ pairwise coprime.  This simple observation is used at four
different scales.  The logical chain is
\begin{equation}\label{eq:roadmap}
 \begin{gathered}
 \text{positive density}
 \Longrightarrow \text{a finite-prime envelope}
 \Longrightarrow \text{a squarefree moving-prefix capacity}\\[1mm]
 \Longrightarrow \text{balanced pair-product cubes}
 \Longrightarrow \text{a cap-set saving}.
 \end{gathered}
\end{equation}
The first two arrows are deterministic.  The difficult point is to build
balanced cubes whose individual word marginals are close to the relevant
product measure even though the extremizing prefix is allowed to move.

For that purpose we use many disjoint prime bands.  On one band, a
five-state coupling creates a common part and three disjoint petals.  If
$B_T$ denotes the event that the three petal weights nearly agree and all
four active labels satisfy lower prefix profiles, then
\[
 \Pp(B_T)\gg w^2,
 \qquad
 \E\bigl[\Pp(B_T\mid S)^2\bigr]\ll w^4.
\]
The first estimate is proved by a whole-band Poisson coupling and two
adaptive anchors.  The second is a fully discrete calculation: two
canonical root pivots give two factors of $w$, and two missing-petal pivots
give two more.  Consequently the conditioned root density has a uniform
$L^2$ bound.  Averaging the active band among many alternatives converts
this $L^2$ control into the $L^1$ flatness needed by the moving-prefix
argument.

The paper is organized in dependency order.  Section~\ref{sec:structure}
proves the elementary pair-product description and records the external
inputs.  Section~\ref{sec:reduction} reduces the theorem to a normalized
squarefree capacity.  Section~\ref{sec:cubes} proves the balanced-cube
transference principle.  Section~\ref{sec:firstmoment} constructs one
prime-band coordinate and proves its first-moment lower bound.
Section~\ref{sec:secondmoment} proves the exact pivot estimates and the
matching second-moment upper bound.  Section~\ref{sec:flattening} flattens
the word marginals by alternative bands, and Section~\ref{sec:completion}
places those bands and completes the proof.  Appendix~\ref{app:verification}
describes the companion verifier and its precise scope.

\section{Structure and external inputs}\label{sec:structure}

\subsection{The pair-product form}

\begin{lemma}[Pair-product form]\label{lem:pair-product}
Three distinct positive integers $a,b,c$ satisfy \eqref{eq:forbidden} if
and only if, after possibly permuting them, there are positive integers
$t,x,y,z$ such that
\begin{equation}\label{eq:pair-product}
 a=txy,\qquad b=txz,\qquad c=tyz,
\end{equation}
and $x,y,z$ are pairwise coprime.
\end{lemma}

\begin{proof}
If \eqref{eq:pair-product} holds, pairwise coprimality gives
\[
 \lcm(a,b)=\lcm(a,c)=\lcm(b,c)=txyz
\]
prime by prime.

Conversely, let $L$ be the common value in \eqref{eq:forbidden}.  Fix a
prime $p$ and write
\[
 A=v_p(a),\qquad B=v_p(b),\qquad C=v_p(c),\qquad M=v_p(L).
\]
The maximum of each pair among $A,B,C$ is $M$.  Therefore at least two of
$A,B,C$ equal $M$.  Define
\[
 x=\frac{L}{c},\qquad y=\frac{L}{b},\qquad z=\frac{L}{a}.
\]
If, for example, $p$ divided both $x$ and $y$, then $C<M$ and $B<M$, which
would contradict $\max(B,C)=M$.  The same argument for the other two pairs
shows that $x,y,z$ are pairwise coprime.

Finally put $t=abc/L^2$.  At the prime $p$ its exponent is
\[
 A+B+C-2M\ge0,
\]
because at least two of $A,B,C$ equal $M$.  Hence $t$ is an integer.  A
direct valuation calculation now gives
\[
 v_p(txy)=A,\qquad v_p(txz)=B,\qquad v_p(tyz)=C.
\]
Since this holds for every prime, \eqref{eq:pair-product} follows.
\end{proof}

For squarefree integers, least common multiple corresponds to union of
prime supports.  Thus Lemma~\ref{lem:pair-product} is the bridge between the
integer problem and the set systems used below.

\subsection{Prime-number estimates}

We use the following standard consequences of the classical zero-free
region and the Brun--Titchmarsh inequality.  The quoted forms follow by
partial summation from the prime number theorem with de la Vall\'ee Poussin
error; see \cite[Chapters~2 and~6]{MontgomeryVaughan} and
\cite[Chapters~5 and~6]{IwaniecKowalski}.  The interval estimate is the
classical Brun--Titchmarsh theorem; see
\cite{MontgomeryVaughanLargeSieve}.

\begin{proposition}[Analytic prime inputs]\label{prop:prime-inputs}
There are constants $B_0,B_1\in\mathbb R$ and $c>0$ such that, for
$x\ge3$,
\begin{align}
 \sum_{p\le x}\frac1p
   &=\log\log x+B_0+O\!\left(\e^{-c\sqrt{\log x}}\right),
                                                        \label{eq:mertens-recip}\\
 \sum_{p\le x}\frac1{p+1}
   &=\log\log x+B_1+O\!\left(\e^{-c\sqrt{\log x}}\right),
                                                        \label{eq:mertens-plus}\\
 \sum_{p\le x}\frac{\log p}{p-1}
   &=(1+o(1))\log x.                                  \label{eq:weighted-prime}
\end{align}
Moreover, whenever $2\le h\le x$,
\begin{equation}\label{eq:brun-titchmarsh}
 \pi(x+h)-\pi(x)\le\frac{2h}{\log h}.
\end{equation}
If $h>x$, the ordinary prime number theorem gives the weaker estimates
needed below.
\end{proposition}

The replacement of $1/p$ by $1/(p+1)$ in
\eqref{eq:mertens-plus} changes only the constant, because
$\sum_p p^{-2}<\infty$ and its tail is $O(1/x)$.

\subsection{The cap-set estimate}

\begin{proposition}[Cap-set bound]\label{prop:cap-set}
Let $A\subseteq\F_3^H$ contain no nonconstant affine line.  Then
\begin{equation}\label{eq:cap-set}
 |A|\le 3\kappacap^H3^H,
 \qquad
 \kappacap:=\frac7{12}2^{2/3}<1.
\end{equation}
\end{proposition}

\begin{proof}
The polynomial-method bound of Ellenberg and Gijswijt
\cite{EllenbergGijswijt} gives
\[
 |A|\le3\min_{0<t<1}
       \bigl((1+t+t^2)t^{-2/3}\bigr)^H.
\]
Taking $t=1/2$ gives
$(1+t+t^2)t^{-2/3}=3\kappacap$.  Also
\[
 \kappacap^3=4\left(\frac7{12}\right)^3
             =\frac{343}{432}<1,
\]
which proves the displayed strict inequality.  The same exact calculation
is reproduced by the companion verifier.
\end{proof}

\section{Reduction to a squarefree moving-prefix capacity}
\label{sec:reduction}

\subsection{The finite-prime envelope}

Let $P$ be a finite set of primes.  Write $\cM_P$ for the multiplicative
monoid generated by $P$, and put
\[
 \delta_P:=\prod_{p\in P}\left(1-\frac1p\right).
\]
Let $b_P(T)$ be the largest size of a safe subset of
$\cM_P\cap[1,T]$, and define
\begin{equation}\label{eq:finite-envelope-capacity}
 C(P):=\delta_P\int_1^\infty b_P(T)\frac{\dd T}{T^2}.
\end{equation}
This integral converges.  Indeed, a member of $\cM_P\cap[1,T]$ is
determined by $|P|$ nonnegative exponents, each $O_P(1+\log T)$, and hence
\begin{equation}\label{eq:smooth-count}
 b_P(T)\le|\cM_P\cap[1,T]|=O_P((1+\log T)^{|P|}).
\end{equation}

\begin{proposition}[Finite-prime envelope]\label{prop:finite-envelope}
For every fixed finite set $P$ of primes,
\begin{equation}\label{eq:finite-envelope}
 \limsup_{N\to\infty}\frac{f(N)}N\le C(P).
\end{equation}
\end{proposition}

\begin{proof}
Fix $P$ throughout the proof.  Every positive integer has a unique
factorization $n=mq$ with $q\in\cM_P$ and
$(m,\prod_{p\in P}p)=1$.  If a safe subset of $[N]$ is restricted to one
fixed $m$-fibre, the resulting set of $q$'s is safe: here $m$ is coprime to
every $q$, so
\[
 \lcm(mq_i,mq_j)=m\lcm(q_i,q_j).
\]
It follows that
\begin{equation}\label{eq:fibre-bound}
 f(N)\le
 \sum_{\substack{m\le N\\(m,\prod_{p\in P}p)=1}}b_P(N/m).
\end{equation}

Finite inclusion--exclusion gives
\[
 A_P(x):=\#\{m\le x:(m,\textstyle\prod_{p\in P}p)=1\}
 =\delta_Px+E_P(x),
 \qquad E_P(x)=O_P(1).
\]
The function $t\mapsto b_P(N/t)$ is nonincreasing, and its total variation
on $[1,N]$ is at most $b_P(N)$.  Stieltjes summation therefore gives
\begin{align}
 \sum_{\substack{m\le N\\(m,\prod_{p\in P}p)=1}}b_P(N/m)
 &=\int_{1^-}^{N}b_P(N/t)\,\dd A_P(t)\notag\\
 &=\delta_P\int_1^N b_P(N/t)\,\dd t+O_P(b_P(N)).
                                                        \label{eq:stieltjes}
\end{align}
For completeness, the error term follows by integration by parts: the two
endpoint terms are $O_P(b_P(N))$, and the remaining integral is bounded by
$\|E_P\|_\infty$ times the total variation of $b_P(N/t)$.

In the main integral set $u=N/t$.  Then
\begin{equation}\label{eq:envelope-integral}
 \delta_P\int_1^N b_P(N/t)\,\dd t
 =\delta_PN\int_1^N b_P(u)\frac{\dd u}{u^2}.
\end{equation}
By \eqref{eq:smooth-count}, $b_P(N)=o(N)$ for fixed $P$.  Divide
\eqref{eq:fibre-bound}--\eqref{eq:envelope-integral} by $N$ and let
$N\to\infty$ to obtain \eqref{eq:finite-envelope}.
\end{proof}

\subsection{Deleting all unit exponents}

For a finite prime set $\cR$ and $S\subseteq\cR$, write
\[
 d(S):=\prod_{p\in S}p,
 \qquad
 Z_{\cR}:=\prod_{p\in\cR}(1+p^{-1}),
 \qquad
 \mu_{\cR}(S):=\frac1{Z_{\cR}d(S)}.
\]
Thus $\mu_{\cR}$ is the product law in which $p$ is selected with
probability $1/(p+1)$.

A family $\mathcal H\subseteq2^{\cR}$ is \emph{admissible} if it contains
no three distinct members $S_1,S_2,S_3$ for which
\begin{equation}\label{eq:equal-unions}
 S_1\cup S_2=S_1\cup S_3=S_2\cup S_3.
\end{equation}
Let $A_{\cR}^{\mathrm{sf}}(x)$ be the largest size of an admissible
subfamily of
\[
 \{S\subseteq\cR:d(S)\le x\},
\]
and adopt the convention
\begin{equation}\label{eq:capacity-zero-convention}
 A_{\cR}^{\mathrm{sf}}(x):=0\qquad(0<x<1).
\end{equation}
Thus $A_{\cR}^{\mathrm{sf}}$ is defined for every $x>0$.  We then define
\begin{equation}\label{eq:squarefree-capacity}
 I_{\cR}^{\mathrm{sf}}
 :=\int_1^\infty A_{\cR}^{\mathrm{sf}}(x)\frac{\dd x}{x^2}.
\end{equation}
The elementary bound
\begin{equation}\label{eq:capacity-trivial}
 0\le\frac{I_{\cR}^{\mathrm{sf}}}{Z_{\cR}}\le1
\end{equation}
follows from
\[
 I_{\cR}^{\mathrm{sf}}
 \le\sum_{S\subseteq\cR}\int_{d(S)}^\infty\frac{\dd x}{x^2}
 =\sum_{S\subseteq\cR}\frac1{d(S)}=Z_{\cR}.
\]

For $E\subseteq P$ and $m\in\cM_P$, set
\[
 U_E(m):=\{p\in E:v_p(m)=1\}.
\]

\begin{proposition}[All-unit deletion]\label{prop:unit-deletion}
For finite prime sets $E\subseteq P$,
\begin{equation}\label{eq:deletion-bound}
 C(P)\le\delta_P
 \sum_{\substack{m\in\cM_P\\U_E(m)=\varnothing}}
 \frac1m I_{E\setminus\supp(m)}^{\mathrm{sf}}.
\end{equation}
Moreover,
\begin{equation}\label{eq:deletion-normalized}
 C(P)\le
 \frac{I_E^{\mathrm{sf}}}{Z_E}+\sum_{p\in E}\frac1{p^2}.
\end{equation}
\end{proposition}

\begin{proof}
Every $q\in\cM_P$ has a unique decomposition
\begin{equation}\label{eq:unit-decomposition}
 q=m\,d(S),
 \qquad
 S=U_E(q),
 \qquad
 U_E(m)=\varnothing,
 \qquad
 S\subseteq E\setminus\supp(m).
\end{equation}
Indeed, $d(S)$ removes exactly one copy of each prime whose exponent in
$q$ is one.  Hence every prime of $E$ has exponent either zero or at least
two in $m$.

Fix $m$ in \eqref{eq:unit-decomposition} and consider a safe family in
$\cM_P\cap[1,T]$.  Its $m$-section is admissible.  Otherwise three supports
satisfying \eqref{eq:equal-unions} would give
\[
 \lcm(md(S_i),md(S_j))=m\,d(S_i\cup S_j)
\]
for each pair, contradicting safety.  Therefore
\[
 b_P(T)\le
 \sum_{\substack{m\in\cM_P\\U_E(m)=\varnothing}}
 A_{E\setminus\supp(m)}^{\mathrm{sf}}(T/m).
\]
Insert this inequality into \eqref{eq:finite-envelope-capacity}.  Positivity
allows termwise integration.  For each eligible $m$, the substitution
$T=mx$ and the convention \eqref{eq:capacity-zero-convention} give
\begin{align*}
 \int_1^\infty
 A_{E\setminus\supp(m)}^{\mathrm{sf}}(T/m)\frac{\dd T}{T^2}
 &=\frac1m\int_{1/m}^\infty
 A_{E\setminus\supp(m)}^{\mathrm{sf}}(x)\frac{\dd x}{x^2}\\
 &=\frac1m\int_1^\infty
 A_{E\setminus\supp(m)}^{\mathrm{sf}}(x)\frac{\dd x}{x^2}
 =\frac1m I_{E\setminus\supp(m)}^{\mathrm{sf}}.
\end{align*}
This proves \eqref{eq:deletion-bound}.

It remains to normalize the right side.  For an eligible $m$, put
$R=E\setminus\supp(m)$ and define
\begin{equation}\label{eq:deletion-law}
 \mathcal W(m):=\delta_P\frac{Z_R}{m}.
\end{equation}
We verify prime by prime that $\mathcal W$ is a probability law.  If
$p\notin E$, summing all valuations gives
\[
 (1-p^{-1})\sum_{k\ge0}p^{-k}=1.
\]
If $p\in E$, valuation zero contributes
\[
 (1-p^{-1})(1+p^{-1})=1-p^{-2},
\]
while all permitted positive valuations, namely $k\ge2$, contribute
\[
 (1-p^{-1})\sum_{k\ge2}p^{-k}=p^{-2}.
\]
Thus the local masses sum to one.  They also show that the events
$\{p\notin R\}$ are independent and satisfy
\begin{equation}\label{eq:deletion-error-law}
 \Pp_{\mathcal W}(p\notin R)=p^{-2}\qquad(p\in E).
\end{equation}

The right side of \eqref{eq:deletion-bound} is exactly
\[
 \E_{\mathcal W}\left[\frac{I_R^{\mathrm{sf}}}{Z_R}\right].
\]
On $\{R=E\}$ the random ratio equals $I_E^{\mathrm{sf}}/Z_E$; off that
event it lies in $[0,1]$ by \eqref{eq:capacity-trivial}.  Hence
\[
 C(P)\le\frac{I_E^{\mathrm{sf}}}{Z_E}
          +\Pp_{\mathcal W}(R\ne E)
 \le\frac{I_E^{\mathrm{sf}}}{Z_E}+\sum_{p\in E}p^{-2},
\]
which is \eqref{eq:deletion-normalized}.
\end{proof}

\begin{corollary}[The remaining target]\label{cor:capacity-target}
Let $P_y=\{p:p\le y\}$.  Suppose that $E_y\subseteq P_y$ satisfies
\begin{equation}\label{eq:capacity-target}
 \min E_y\longrightarrow\infty,
 \qquad
 \frac{I_{E_y}^{\mathrm{sf}}}{Z_{E_y}}\longrightarrow0.
\end{equation}
Then $f(N)=o(N)$.
\end{corollary}

\begin{proof}
Apply Proposition~\ref{prop:unit-deletion} with $P=P_y$ and $E=E_y$.
Since $\min E_y\to\infty$,
\[
 \sum_{p\in E_y}p^{-2}
 \le\sum_{n\ge\min E_y}n^{-2}\longrightarrow0.
\]
Thus $C(P_y)\to0$.  Proposition~\ref{prop:finite-envelope} gives, for each
fixed $y$,
\[
 \limsup_{N\to\infty}\frac{f(N)}N\le C(P_y).
\]
Letting $y\to\infty$ proves the corollary.
\end{proof}

\section{Balanced cubes and moving prefixes}\label{sec:cubes}

\subsection{Pair-product cubes}

A dimension-$H$ \emph{pair-product cube} consists of a common prime support
and, for each coordinate $1\le i\le H$, three nonempty prime supports
$x_i,y_i,z_i$.  All $3H$ petals are mutually disjoint, and the common
support is disjoint from every petal.  A word
$\omega=(\omega_1,\ldots,\omega_H)\in\F_3^H$ takes one of the following
three contributions in coordinate $i$:
\begin{center}
\begin{tabular}{c c c}
\toprule
state $\omega_i$ & chosen petals & omitted petal\\
\midrule
$0$ & $y_i\cup z_i$ & $x_i$\\
$1$ & $x_i\cup z_i$ & $y_i$\\
$2$ & $x_i\cup y_i$ & $z_i$\\
\bottomrule
\end{tabular}
\end{center}
The support of the word is the union of its $H$ coordinate contributions
and the common support.

\begin{lemma}[Cap-set saving on one cube]\label{lem:cube-saving}
An admissible family occupies at most a proportion
$3\kappacap^H$ of the $3^H$ words of a pair-product cube.
\end{lemma}

\begin{proof}
The word map is injective: if two words first differ in coordinate $i$,
then one of the nonempty petals $x_i,y_i,z_i$ occurs in exactly one of the
two supports.

Consider a nonconstant affine line
\[
 \omega,\qquad \omega+v,\qquad \omega+2v
 \quad\text{in }\F_3^H.
\]
In a coordinate where $v_i=0$, all three words make the same choice.  In a
coordinate where $v_i\ne0$, the three states are $0,1,2$ in some order.
The pairwise union of any two corresponding coordinate contributions is
then $x_i\cup y_i\cup z_i$.  It follows coordinate by coordinate that the
three full supports have identical pairwise unions.  Consequently the
preimage of an admissible family contains no nonconstant affine line.
Proposition~\ref{prop:cap-set} bounds that preimage by
$3\kappacap^H3^H$, proving the lemma.
\end{proof}

\subsection{A joint law for a moving prefix}

The extremal family in the definition of
$A_{\cR}^{\mathrm{sf}}(x)$ depends on $x$.  It is therefore not enough to
construct a good cube at one fixed cutoff.  The next proposition couples
the cube and the cutoff so that all words are tested against the same
moving prefix.

\begin{proposition}[Joint-prefix transference]\label{prop:joint-prefix}
Let $\cR$ be finite, and let $\lambda$ be a probability law on
dimension-$H$ pair-product cubes in $2^{\cR}$.  For a word $\omega$, let
$\nu_\omega$ be its support marginal.  Suppose that
\begin{equation}\label{eq:word-density}
 \frac{\dd\nu_\omega}{\dd\mu_{\cR}}=G_\omega,
 \qquad
 \E_{\mu_{\cR}}|G_\omega-1|\le\eps
\end{equation}
for every $\omega\in\F_3^H$.  Suppose also that $0\le\zeta\le1$ and, for
every sampled cube $c$ and every word $\omega$,
\begin{equation}\label{eq:log-balance}
 \left|\log d(S_\omega(c))-\log\mathcal G(c)\right|\le\zeta,
 \qquad
 \mathcal G(c):=
 \left(\prod_{\tau\in\F_3^H}d(S_\tau(c))\right)^{1/3^H}.
\end{equation}
Then
\begin{equation}\label{eq:joint-prefix-bound}
 \frac{I_{\cR}^{\mathrm{sf}}}{Z_{\cR}}
 \le 3\kappacap^H+O(\eps+\zeta),
\end{equation}
where the implied constant is absolute.
\end{proposition}

\begin{proof}
We separate the construction into four steps.

\smallskip
\noindent\emph{Step 1: a joint cube--cutoff measure.}
For a cube $c$, put
\[
 D(c):=\max_{\omega\in\F_3^H}d(S_\omega(c)).
\]
Define a finite measure on the cube space times $[1,\infty)$ by
\begin{equation}\label{eq:joint-measure}
 \dd\Lambda(c,x):=
 \mathcal G(c)\,\1_{\{x\ge D(c)\}}
 \frac{\dd x}{x^2}\,\dd\lambda(c).
\end{equation}
Since $D(c)\ge\mathcal G(c)$, while \eqref{eq:log-balance} gives
$D(c)\le\e^\zeta\mathcal G(c)$,
\begin{equation}\label{eq:joint-mass}
 \e^{-\zeta}\le
 \|\Lambda\|
 =\E_\lambda\frac{\mathcal G(c)}{D(c)}
 \le1.
\end{equation}

\smallskip
\noindent\emph{Step 2: the weighted mass of one fibre.}
Fix a word $\omega$ and a support $S\subseteq\cR$.  Whenever
$S_\omega(c)=S$, the balance assumption implies
\begin{equation}\label{eq:fibre-balance}
 \e^{-\zeta}d(S)\le\mathcal G(c)\le\e^\zeta d(S),
 \qquad
 D(c)\le\e^{2\zeta}d(S).
\end{equation}
The second inequality follows because every word product is at most
$\e^\zeta\mathcal G(c)$ and
$\mathcal G(c)\le\e^\zeta d(S)$.  By \eqref{eq:word-density},
\[
 \nu_\omega(S)=\frac{G_\omega(S)}{Z_{\cR}d(S)}.
\]
Consequently the weighted fibre mass
\[
 M_\omega(S):=
 \int_{\{c:S_\omega(c)=S\}}\mathcal G(c)\,\dd\lambda(c)
\]
satisfies the pointwise bounds
\begin{equation}\label{eq:weighted-fibre}
 \frac{\e^{-\zeta}G_\omega(S)}{Z_{\cR}}
 \le M_\omega(S)\le
 \frac{\e^\zeta G_\omega(S)}{Z_{\cR}}.
\end{equation}

\smallskip
\noindent\emph{Step 3: comparison with the canonical prefix law.}
We use the following convention.  If $\sigma$ is a finite signed measure,
then $\|\sigma\|_{\mathrm{var}}:=|\sigma|(\Omega)$ denotes its total
variation norm.  For probability measures $P,Q$ we write
\[
 d_{\TV}(P,Q):=\sup_E|P(E)-Q(E)|
 =\frac12\|P-Q\|_{\mathrm{var}}.
\]
Let $\Lambda_\omega$ be the pushforward of $\Lambda$ under
$(c,x)\mapsto(S_\omega(c),x)$.  With respect to
\[
 \rho:=\sum_{S\subseteq\cR}\delta_S\otimes\frac{\dd x}{x^2}
 \quad\text{on }2^{\cR}\times[1,\infty),
\]
its density is
\[
 L_\omega(S,x):=
 \int_{\{c:S_\omega(c)=S\}}
 \mathcal G(c)\1_{\{D(c)\le x\}}\,\dd\lambda(c).
\]
Let $\Pi$ be the probability measure whose $\rho$-density is
\begin{equation}\label{eq:canonical-prefix-law}
 \frac{\dd\Pi}{\dd\rho}(S,x)
 :=\frac1{Z_{\cR}}\1_{\{x\ge d(S)\}}.
\end{equation}
Indeed,
\[
 \Pi(2^{\cR}\times[1,\infty))
 =\frac1{Z_{\cR}}\sum_{S\subseteq\cR}\frac1{d(S)}=1.
\]

Fix $S\subseteq\cR$ and abbreviate
$G=G_\omega(S)$ and $M=M_\omega(S)$.  If $G>0$,
\eqref{eq:weighted-fibre} permits us to write
\[
 M=\frac{\theta_SG}{Z_{\cR}},
 \qquad \e^{-\zeta}\le\theta_S\le\e^\zeta.
\]
If $G=0$, then \eqref{eq:weighted-fibre} gives $M=0$, and the same formula
holds after setting $\theta_S=1$.  Since $0\le\zeta\le1$,
\[
 |\theta_S-1|\le\e^\zeta-1\le2\zeta.
\]
On the fibre $S_\omega(c)=S$, one has $D(c)\ge d(S)$ and, by
\eqref{eq:fibre-balance}, $D(c)\le\e^{2\zeta}d(S)$.  Consequently
$L_\omega(S,x)=0$ for $x<d(S)$ and
$L_\omega(S,x)=M_\omega(S)$ for $x\ge\e^{2\zeta}d(S)$.  On the latter
region,
\begin{align}
 &\sum_{S\subseteq\cR}
 \int_{\e^{2\zeta}d(S)}^\infty
 \left|L_\omega(S,x)-\frac1{Z_{\cR}}\right|
 \frac{\dd x}{x^2}\notag\\
 &=\e^{-2\zeta}\sum_{S\subseteq\cR}
 \frac{|M_\omega(S)-Z_{\cR}^{-1}|}{d(S)}\notag\\
 &\le
 \sum_{S\subseteq\cR}
 \frac{|G_\omega(S)-1|+2\zeta G_\omega(S)}
      {Z_{\cR}d(S)}
 \le\eps+2\zeta.                                  \label{eq:high-region}
\end{align}
Here the last line uses
$\E_{\mu_{\cR}}|G_\omega-1|\le\eps$ and
$\E_{\mu_{\cR}}G_\omega=1$.

It remains to control the transition strip
\[
 \mathcal T:=\{(S,x):d(S)\le x<\e^{2\zeta}d(S)\}.
\]
Its $\Pi$-mass is exactly
\begin{equation}\label{eq:strip-pi}
 \Pi(\mathcal T)
 =(1-\e^{-2\zeta})
 \sum_{S\subseteq\cR}\frac1{Z_{\cR}d(S)}
 =1-\e^{-2\zeta}\le2\zeta.
\end{equation}
For the other marginal, Tonelli's theorem and nonnegativity give
\begin{align}
 \Lambda_\omega(\mathcal T)
 &=\int \mathcal G(c)
   \int_{d(S_\omega(c))}^{\e^{2\zeta}d(S_\omega(c))}
   \1_{\{D(c)\le x\}}\frac{\dd x}{x^2}\,\dd\lambda(c)\notag\\
 &\le(1-\e^{-2\zeta})
   \int\frac{\mathcal G(c)}{d(S_\omega(c))}\,\dd\lambda(c)\notag\\
 &\le\e^\zeta(1-\e^{-2\zeta})
 \le6\zeta.                                      \label{eq:strip-lambda}
\end{align}
The penultimate inequality follows from \eqref{eq:log-balance}; the last
uses $\e^\zeta\le\e<3$ and
$1-\e^{-2\zeta}\le2\zeta$.  On $\mathcal T$ the variation norm of the
difference is at most the sum of the two masses, while below $x=d(S)$
both measures vanish.  Combining
\eqref{eq:high-region}--\eqref{eq:strip-lambda} yields
\begin{equation}\label{eq:unnormalized-prefix-TV}
 \|\Lambda_\omega-\Pi\|_{\mathrm{var}}
 \le\eps+10\zeta.
\end{equation}

Put $a_\Lambda:=\|\Lambda\|$.  Every $\Lambda_\omega$ has mass
$a_\Lambda$, and \eqref{eq:joint-mass} gives
$\e^{-\zeta}\le a_\Lambda\le1$.  The word marginal
$\overline\Lambda_\omega$ of
$\overline\Lambda:=a_\Lambda^{-1}\Lambda$ is
$a_\Lambda^{-1}\Lambda_\omega$.  Therefore
\begin{align}
 \|\overline\Lambda_\omega-\Pi\|_{\mathrm{var}}
 &\le
 \|a_\Lambda^{-1}\Lambda_\omega-\Lambda_\omega\|_{\mathrm{var}}
 +\|\Lambda_\omega-\Pi\|_{\mathrm{var}}\notag\\
 &=(1-a_\Lambda)+\|\Lambda_\omega-\Pi\|_{\mathrm{var}}
 \le\eps+11\zeta,                              \label{eq:normalized-prefix-TV}
\end{align}
because $1-a_\Lambda\le1-\e^{-\zeta}\le\zeta$.  Thus, uniformly in
$\omega$,
\begin{equation}\label{eq:normalized-prefix-distance}
 d_{\TV}(\overline\Lambda_\omega,\Pi)
 \le\frac{\eps+11\zeta}{2}.
\end{equation}

\smallskip
\noindent\emph{Step 4: the moving extremizer.}
Fix a lexicographic order on the finite set $2^{\cR}$, and use the induced
lexicographic order on its subfamilies.  For each $x\ge1$, choose the first
admissible family
$\mathcal F_x\subseteq\{S:d(S)\le x\}$ of maximum size.  Only finitely many
prefixes occur, so $x\mapsto\mathcal F_x$ is piecewise constant and hence
measurable.  Set
\[
 \mathcal E:=\{(S,x):S\in\mathcal F_x\}.
\]
Because $\mathcal F_x$ is contained in the $x$-prefix, Tonelli's theorem
gives
\begin{align}
 \Pi(\mathcal E)
 &=\frac1{Z_{\cR}}
   \sum_{S\subseteq\cR}\int_{d(S)}^\infty
      \1_{\{S\in\mathcal F_x\}}\frac{\dd x}{x^2}\notag\\
 &=\frac1{Z_{\cR}}\int_1^\infty
      |\mathcal F_x|\frac{\dd x}{x^2}
 =\frac{I_{\cR}^{\mathrm{sf}}}{Z_{\cR}}.       \label{eq:prefix-occupancy}
\end{align}

For $\overline\Lambda$-almost every $(c,x)$, one has
$x\ge D(c)$, so all $3^H$ word supports lie in the same $x$-prefix.
Since $\mathcal F_x$ is admissible, Lemma~\ref{lem:cube-saving} implies
\[
 \frac1{3^H}\sum_{\omega\in\F_3^H}
 \1_{\{S_\omega(c)\in\mathcal F_x\}}
 \le3\kappacap^H.
\]
Integrating this pointwise inequality and using that
$\overline\Lambda_\omega$ is the $(S_\omega,x)$-pushforward of
$\overline\Lambda$, we obtain
\begin{align}
 \frac1{3^H}\sum_{\omega\in\F_3^H}
 \overline\Lambda_\omega(\mathcal E)
 &=\int\frac1{3^H}\sum_{\omega\in\F_3^H}
   \1_{\{S_\omega(c)\in\mathcal F_x\}}
   \,\dd\overline\Lambda(c,x)\notag\\
 &\le3\kappacap^H.                              \label{eq:word-average-upper}
\end{align}
On the other hand, \eqref{eq:normalized-prefix-distance} and
\eqref{eq:prefix-occupancy} imply, for every $\omega$,
\[
 \overline\Lambda_\omega(\mathcal E)
 \ge\frac{I_{\cR}^{\mathrm{sf}}}{Z_{\cR}}
      -\frac{\eps+11\zeta}{2}.
\]
Averaging this inequality over all words and comparing with
\eqref{eq:word-average-upper} gives
\[
 \frac{I_{\cR}^{\mathrm{sf}}}{Z_{\cR}}
 \le3\kappacap^H+\frac{\eps+11\zeta}{2},
\]
which proves \eqref{eq:joint-prefix-bound}.
\end{proof}

\section{One prime band: the model and its first moment}
\label{sec:firstmoment}

\subsection{Explicit parameters and normalized prime weights}

We make all structural constants explicit:
\begin{equation}\label{eq:parameters}
 a=\frac{23}{25},\qquad
 a_0=\frac{24}{25},\qquad
 \vartheta=\frac1{200},\qquad
 \alpha=\frac8{25},\qquad
 q:=3\alpha=\frac{24}{25}.
\end{equation}
These choices satisfy
\begin{equation}\label{eq:parameter-inequalities}
 \frac1{\log3}<a<a_0<1,
 \qquad
 a(1-\vartheta)\log3>1.
\end{equation}
Here is an exact certificate.  The positive series
\[
 \log3
 =2\sum_{n\ge0}\frac{(1/2)^{2n+1}}{2n+1}
\]
gives, on retaining its first five terms,
\[
 \log3>
 1+\frac1{12}+\frac1{80}+\frac1{448}+\frac1{2304}
 >\frac{549}{500}.
\]
Consequently
\[
 a(1-\vartheta)\log3
 >
 \frac{23}{25}\frac{199}{200}\frac{549}{500}
 =\frac{2512773}{2500000}>1.
\]
This also implies $a\log3>1$.  The companion verifier reproduces these
comparisons with exact rational arithmetic.

Let $T\to\infty$, and define
\begin{equation}\label{eq:band-definitions}
 \cB(T):=\{p:\e^{T^\vartheta}<p\le\e^T\},
 \qquad
 u_p:=\frac{\log p}{T},
 \qquad
 s_p:=-\log u_p.
\end{equation}
Thus $u_p\in(T^{\vartheta-1},1]$.  The largest possible depth is
\begin{equation}\label{eq:horizon}
 \lambda_T:=(1-\vartheta)\log T.
\end{equation}
Let $\eta=\eta(T)$ satisfy
\begin{equation}\label{eq:slow-eta}
 \eta\longrightarrow0,
 \qquad
 \log(1/\eta)=o(\log T),
 \qquad
 w:=\frac{\eta}{T}.
\end{equation}
Equivalently, $\eta=T^{-o(1)}$.  All statements in this section hold for
every such function $\eta$ once $T$ is sufficiently large; the threshold
may depend on the function, but all displayed constants are uniform in
$T$ and $\eta$.

\subsection{Reciprocal-prime mass on a normalized band}

The upper bound below is expressed in terms of $w$, not the actual length
of the interval.  This is important because reciprocal-prime measure is
atomic.

\begin{lemma}[Local reciprocal-prime mass]\label{lem:local-prime-mass}
Fix $C_0>0$.  For every sufficiently large $T$, every interval
$I\subseteq[T^{\vartheta-1},1]$ of length at most $C_0w$ satisfies
\begin{equation}\label{eq:local-upper}
 \sum_{\substack{p\in\cB(T)\\u_p\in I}}\frac1p
 \ll_{C_0}\frac{w}{\inf I}.
\end{equation}
This includes degenerate intervals and intervals containing only one prime.

If $0<r_0<r_1<1$ and $c_0>0$ are fixed, then, uniformly for
$t\in[r_0,r_1]$,
\begin{equation}\label{eq:local-lower}
 \sum_{\substack{p\in\cB(T)\\t\le u_p\le t+c_0w}}\frac1p
 =\log\!\left(1+\frac{c_0w}{t}\right)+o(w)
 =(1+o(1))\frac{c_0w}{t}.
\end{equation}
Both statements remain valid with $1/p$ replaced by $1/(p+1)$.
\end{lemma}

\begin{proof}
Put $\epsilon_T=T^{\vartheta-1}$.  Since
\[
 \frac{w}{\epsilon_T}=\frac{\eta}{T^\vartheta}=o(1),
\]
we may enlarge $I$, only in an available direction at a band endpoint, to
an interval $[r,s]\subseteq[\epsilon_T,1]$ such that
\begin{equation}\label{eq:interval-enlargement}
 I\subseteq[r,s],\qquad
 w\le s-r\le(C_0+2)w,\qquad
 r\ge\frac12\inf I.
\end{equation}
Set $x=\e^{Tr}$ and $\Delta=\e^{Ts}-\e^{Tr}$.  Because
$T(s-r)\asymp_{C_0}\eta$,
\[
 \Delta\asymp_{C_0}x\eta,\qquad
 \Delta\le x,\qquad
 \log\Delta=Tr+\log\eta+O_{C_0}(1).
\]
Now $Tr\ge T^\vartheta/2$, while
$|\log\eta|=o(\log T)=o(T^\vartheta)$.  Hence, for large $T$,
\[
 \log\Delta\ge\frac14T\inf I
 \quad\text{and}\quad \Delta\ge2.
\]
Brun--Titchmarsh, Proposition~\ref{prop:prime-inputs}, gives
\[
 \pi(x+\Delta)-\pi(x)
 \ll_{C_0}\frac{x\eta}{T\inf I}.
\]
Allowing one possible prime at a closed lower endpoint,
\begin{align*}
 \sum_{\substack{p\in\cB(T)\\u_p\in I}}\frac1p
 &\le\frac{\pi(x+\Delta)-\pi(x)+1}{x}\\
 &\ll_{C_0}\frac{\eta}{T\inf I}+\e^{-T^\vartheta}
 \ll_{C_0}\frac{w}{\inf I}.
\end{align*}
 The last absorption uses $\e^{-T^\vartheta}=o(w)$.  The enlargement in
 \eqref{eq:interval-enlargement} was made for every $I$, including a
 degenerate interval.  Equivalently, a degenerate $I$ contains at most one
 normalized prime weight, whose reciprocal contribution is at most
 $\e^{-T^\vartheta}=o(w)$; thus single-atom intervals cause no exception.

For \eqref{eq:local-lower}, subtract
\eqref{eq:mertens-recip} at $\e^{Tt}$ and
$\e^{T(t+c_0w)}$.  Uniformly for $t\in[r_0,r_1]$ this gives
\[
 \sum_{\substack{p\in\cB(T)\\t\le u_p\le t+c_0w}}\frac1p
 =
 \log\frac{t+c_0w}{t}
 +O(\e^{-c\sqrt{Tr_0}})+O(\e^{-Tr_0}).
\]
Both error terms are $o(w)$ because $\eta=T^{-o(1)}$.  Finally,
\[
 0\le
 \sum_{p\in\cB(T)}
 \left(\frac1p-\frac1{p+1}\right)
 \le\sum_{n>\e^{T^\vartheta}}\frac1{n^2}
 \ll\e^{-T^\vartheta}=o(w),
\]
which proves the versions with $1/(p+1)$.
\end{proof}

\subsection{The five-state coupling and the target event}

For $p\in\cB(T)$ put
\[
 r_p:=\frac1{p+1},
 \qquad
 h_p:=\frac{r_p}{3}=\frac1{3(p+1)}.
\]
Independently at each prime, choose one of the five labels
\[
 \varnothing,\quad c,\quad x,\quad y,\quad z
\]
with probabilities
\begin{equation}\label{eq:five-state-law}
 \Pp(c)=\Pp(x)=\Pp(y)=\Pp(z)=h_p,
 \qquad
 \Pp(\varnothing)=1-4h_p.
\end{equation}
These probabilities are nonnegative because $r_p\le1/3$.  From one labelled
configuration form the following three state supports:
\begin{center}
\begin{tabular}{c c c}
\toprule
state & support on the band & missing petal\\
\midrule
$S^{(x)}$ & $c\cup y\cup z$ & $x$\\
$S^{(y)}$ & $c\cup x\cup z$ & $y$\\
$S^{(z)}$ & $c\cup x\cup y$ & $z$\\
\bottomrule
\end{tabular}
\end{center}
A fixed prime belongs to any one state with probability $3h_p=r_p$.
Independence across primes therefore shows that each state has exactly the
product law $\mu_{\cB(T)}$.

For a labelled configuration define the normalized petal weights
\begin{equation}\label{eq:petal-weights}
 X:=\sum_{p:\,x(p)}u_p,\qquad
 Y:=\sum_{p:\,y(p)}u_p,\qquad
 Z:=\sum_{p:\,z(p)}u_p
\end{equation}
and the four prefix counts
\begin{equation}\label{eq:label-counts}
 N_\ell(s):=\#\{p:\ell(p)=\ell,\ s_p\le s\},
 \qquad \ell\in\{c,x,y,z\}.
\end{equation}

We now fix, once and for all,
\begin{equation}\label{eq:fixed-window-depth}
 J:=\left[\frac9{20},\frac{11}{20}\right],
 \qquad R_0:=75.
\end{equation}
The parameter choice gives, for every real $s\ge R_0$,
\begin{equation}\label{eq:profile-dominates}
 3\lfloor\alpha s\rfloor
 \ge3\alpha s-3=a_0s-3\ge as,
\end{equation}
because $(a_0-a)R_0=3$.  Also
$\lfloor\alpha R_0\rfloor=24\ge1$.

Define the symmetric event
\begin{align}
 B_T:=\{&
 X,Y,Z\in J,\quad
 \max(X,Y,Z)-\min(X,Y,Z)\le w;\notag\\
 &N_\ell(s)\ge\lfloor\alpha s\rfloor
 \text{ for }\ell\in\{c,x,y,z\}
 \text{ and }R_0\le s\le\lambda_T\}.
                                                        \label{eq:target-event}
\end{align}
Fix the first state $S=S^{(x)}$ and set
\begin{equation}\label{eq:rooted-functions}
 g_T(S):=\Pp(B_T\mid S),
 \qquad
 \beta_T:=\Pp(B_T).
\end{equation}
The event $B_T$ is invariant under every permutation of $x,y,z$.
Consequently the same function $g_T$ describes all three states.  Bayes'
formula gives the exact Radon--Nikodym identity
\begin{equation}\label{eq:bayes-density}
 \frac{\dd\nu_T}{\dd\mu_{\cB(T)}}(S)
 =\frac{g_T(S)}{\beta_T},
\end{equation}
where $\nu_T$ is the law of any state conditioned on $B_T$.

The next three lemmas prepare the lower bound
$\beta_T\gg w^2$.

\begin{lemma}[Whole-band Poisson coupling]\label{lem:poisson-coupling}
Let $\mathsf Q_T$ be the categorical law
\eqref{eq:five-state-law}.  Regard it as a law on four nonnegative integer
arrays, one for each active label.  Let $\widetilde{\mathsf Q}_T$ be the law
of independent random variables
\[
 K_{\ell,p}\sim\operatorname{Poisson}(h_p),
 \qquad
 \ell\in\{c,x,y,z\},\quad p\in\cB(T).
\]
Then
\begin{equation}\label{eq:whole-band-tv}
 d_{\TV}(\mathsf Q_T,\widetilde{\mathsf Q}_T)
 \le16\sum_{p\in\cB(T)}h_p^2
 \ll\e^{-T^\vartheta}
 =o(w^A)
\end{equation}
for every fixed $A>0$.
\end{lemma}

\begin{proof}
Fix $p$ and write $h=h_p$.  Under the Poisson law, the total number of
points at $p$ is Poisson with mean $4h$.  The empty configuration and each
of the four singleton configurations have probabilities
$\e^{-4h}$ and $h\e^{-4h}$, respectively.  Therefore
\begin{align*}
 2d_{\TV}(\mathsf Q_p,\widetilde{\mathsf Q}_p)
 &\le |1-4h-\e^{-4h}|
      +4h(1-\e^{-4h})
      +\Pp(\operatorname{Poisson}(4h)\ge2).
\end{align*}
The elementary inequalities
\[
 1-x\le\e^{-x}\le1-x+\frac{x^2}{2},
 \qquad
 1-\e^{-x}\le x
 \quad(x\ge0)
\]
bound the first two terms by $8h^2$ and $16h^2$.  If
$N\sim\operatorname{Poisson}(4h)$, then
\[
 \Pp(N\ge2)\le\frac{\E[N(N-1)]}{2}=8h^2.
\]
Thus the local total-variation distance is at most $16h^2$.

Couple the arrays prime by prime.  The probability that any local coupling
fails is at most the sum of the local distances, which proves the first
inequality in \eqref{eq:whole-band-tv}.  Since $h_p\le1/(3p)$,
\[
 \sum_{p\in\cB(T)}h_p^2
 \le\frac19\sum_{n>\e^{T^\vartheta}}\frac1{n^2}
 \ll\e^{-T^\vartheta}.
\]
Finally,
\[
 \log(1/w)=\log T+\log(1/\eta)=(1+o(1))\log T,
\]
whereas $T^\vartheta/\log T\to\infty$.  Hence
$\e^{-T^\vartheta}=o(w^A)$ for every fixed $A$.
\end{proof}

On the common array space define
\begin{equation}\label{eq:poisson-observables}
 X_\ell^{\mathrm P}:=\sum_{p\in\cB(T)}u_pK_{\ell,p},
 \qquad
 N_\ell^{\mathrm P}(s):=
 \sum_{\substack{p\in\cB(T)\\s_p\le s}}K_{\ell,p}.
\end{equation}
Let $B_T^{\mathrm P}$ be the event obtained from
\eqref{eq:target-event} by replacing $X,Y,Z$ with
$X_x^{\mathrm P},X_y^{\mathrm P},X_z^{\mathrm P}$ and every $N_\ell$
with $N_\ell^{\mathrm P}$.  Multiplicities are counted under the Poisson
law.  On the categorical subspace, where each entry is zero or one and at
most one label is nonzero at a prime, $B_T^{\mathrm P}$ is literally
$B_T$.  Thus
\begin{equation}\label{eq:event-identification}
 \mathsf Q_T(B_T^{\mathrm P})
 =\mathsf Q_T(B_T)=\beta_T.
\end{equation}

For a depth interval $D$, write
\[
 \Lambda_T(D):=
 \sum_{\substack{p\in\cB(T)\\s_p\in D}}h_p.
\]

\begin{lemma}[Uniform prime time change and small tails]
\label{lem:prime-time-change}
Uniformly for $0\le t\le\lambda_T-R_0$,
\begin{equation}\label{eq:time-change}
 \Lambda_T((R_0,R_0+t])
 =\frac t3+\varepsilon_T(t),
 \qquad
 \sup_t|\varepsilon_T(t)|=o(1).
\end{equation}
For $\ell\in\{c,x,y,z\}$, put
\begin{equation}\label{eq:tail-weight}
 V_{\ell,T}:=
 \sum_{\substack{p\in\cB(T)\\s_p>R_0}}u_pK_{\ell,p}.
\end{equation}
For all sufficiently large $T$,
\begin{equation}\label{eq:tail-mean}
 \E_{\widetilde{\mathsf Q}_T}V_{\ell,T}\le\e^{-R_0}.
\end{equation}
Finally, for every fixed interval $D\subseteq[0,R_0]$,
\begin{equation}\label{eq:fixed-depth-intensity}
 \Lambda_T(D)\longrightarrow\frac{|D|}{3}.
\end{equation}
All estimates hold simultaneously for the four labels.
\end{lemma}

\begin{proof}
Let
\[
 H_+(x):=\sum_{p\le x}\frac1{p+1}.
\]
For $0\le t\le\lambda_T-R_0$,
\begin{align}
 3\Lambda_T((R_0,R_0+t])
 &=
 H_+(\e^{T\e^{-R_0}})
 -H_+(\e^{T\e^{-(R_0+t)}})
 +O(\e^{-T^\vartheta}).                    \label{eq:time-change-difference}
\end{align}
The logarithms of both prime cutoffs are at least $T^\vartheta$.  By
\eqref{eq:mertens-plus}, the error in
\eqref{eq:time-change-difference} is
$O(\e^{-cT^{\vartheta/2}})$ uniformly in $t$.  Since
\[
 \log\log(\e^{T\e^{-s}})=\log T-s,
\]
the main term in \eqref{eq:time-change-difference} is $t$.  This proves
\eqref{eq:time-change}.  The same subtraction over a fixed depth interval
proves \eqref{eq:fixed-depth-intensity}.

For the weighted tail, positivity and
\eqref{eq:weighted-prime} give
\begin{align*}
 \E V_{\ell,T}
 &=\frac1{3T}
 \sum_{\substack{\e^{T^\vartheta}<p<\e^{T\e^{-R_0}}}}
 \frac{\log p}{p+1}\\
 &\le\frac1{3T}
 \sum_{p\le\e^{T\e^{-R_0}}}\frac{\log p}{p-1}
 =\left(\frac13+o(1)\right)\e^{-R_0}.
\end{align*}
The last expression is at most $\e^{-R_0}$ for all sufficiently large
$T$, proving \eqref{eq:tail-mean}.
\end{proof}

\begin{lemma}[A uniform Poisson lower profile]\label{lem:poisson-survival}
Let $\Pi(v)$ be a rate-one Poisson process.  With
\[
 q=\frac{24}{25},
 \qquad
 \gamma:=\frac1{50},
\]
one has, for every $b\ge0$,
\begin{equation}\label{eq:survival}
 \Pp\bigl(\Pi(v)+b\ge qv\text{ for every }v\ge0\bigr)
 \ge1-\e^{-\gamma b}.
\end{equation}
\end{lemma}

\begin{proof}
The inequality $\e^{-x}\le1-x+x^2/2$ gives
\[
 q\gamma+\e^{-\gamma}-1
 \le(q-1)\gamma+\frac{\gamma^2}{2}
 =-\frac3{5000}<0.
\]
It follows from independent increments that
\[
 M(v):=\exp\!\bigl(\gamma(qv-\Pi(v))\bigr)
\]
is a nonnegative supermartingale.

We include the short maximal argument, including the harmless strictness at
the boundary.  Fix $L<\infty$ and $\epsilon>0$.  On a finite time grid in
$[0,L]$, stop $M$ at the first grid point at which it reaches
$\e^{\gamma(b+\epsilon)}$.  Conditioning on whether the stopping time has
already occurred shows directly that the stopped process is still a
nonnegative supermartingale.  Its expectation is therefore at most
$M(0)=1$, so the grid crossing probability is at most
$\e^{-\gamma(b+\epsilon)}$.  Apply this to nested dyadic grids.  Right
continuity detects every strict crossing, and hence
\[
 \Pp\bigl(\exists v\in[0,L]:qv-\Pi(v)>b+\epsilon\bigr)
 \le\e^{-\gamma(b+\epsilon)}.
\]
Let first $\epsilon\downarrow0$ and then $L\to\infty$.  We obtain
\[
 \Pp\bigl(\exists v\ge0:qv-\Pi(v)>b\bigr)\le\e^{-\gamma b}.
\]
The complement is exactly the event in \eqref{eq:survival}.
\end{proof}

\begin{proposition}[Profiled anchor lower bound]\label{prop:first-moment}
There is an absolute constant $c>0$ such that, for every function
$\eta(T)$ satisfying \eqref{eq:slow-eta} and all sufficiently large $T$,
\begin{equation}\label{eq:first-moment}
 \beta_T=\Pp(B_T)\ge c w^2.
\end{equation}
The constant $c$ is independent of $T$ and $\eta$.
\end{proposition}

\begin{proof}
We construct under $\widetilde{\mathsf Q}_T$ an event contained in
$B_T^{\mathrm P}$ whose probability is at least a fixed multiple of $w^2$.
The whole-band coupling will then transfer it to the categorical law.

\smallskip
\noindent\emph{Step 1: fixed numerical choices.}
Set
\begin{equation}\label{eq:first-moment-constants}
 \delta:=\frac1{100},
 \qquad
 B:=350,
 \qquad
 m:=\lceil\alpha R_0\rceil+B=374.
\end{equation}
These constants satisfy the three inequalities needed below:
\begin{equation}\label{eq:constant-certificate}
 4\e^{-\gamma B/2}<\frac18,
 \qquad
 m\e^{1-R_0}<\frac{\delta}{2},
 \qquad
 \frac{8\e^{-R_0}}{\delta}<\frac18.
\end{equation}
We record exact witnesses.  First,
\[
 \e^{7/2}>
 \sum_{n=0}^{7}\frac{(7/2)^n}{n!}
 =\frac{66007}{2048}>32,
\]
which proves the first inequality.  Next,
\[
 374\e^{-74}
 <\frac{374\cdot4!}{74^4}
 =\frac{561}{1874161}<\frac1{200}.
\]
Finally,
\[
 800\e^{-75}
 <\frac{800\cdot3!}{75^3}
 =\frac{64}{5625}<\frac18.
\]
Only positive terms of the exponential series were used.  These
calculations are also checked exactly in the companion verifier.

\smallskip
\noindent\emph{Step 2: finite-depth buffers.}
Let
\[
 \mathcal K:=\left[\frac9{20},\frac{11}{20}\right),
 \qquad
 \mathcal A:=\left[\frac{12}{25},\frac{13}{25}\right).
\]
For a half-open weight interval $\mathcal I$, write
$D_{\mathcal I}:=\{-\log u:u\in\mathcal I\}$; its image is again taken
with the corresponding half-open endpoint convention.  Choose the four
explicit, pairwise disjoint depth intervals
\begin{align*}
 Q_c&=(R_0-1,R_0-7/8],&
 Q_x&=(R_0-3/4,R_0-5/8],\\
 Q_y&=(R_0-1/2,R_0-3/8],&
 Q_z&=(R_0-1/4,R_0-1/8].
\end{align*}
All lie in $(R_0-1,R_0)$.

Let $F_{\mathrm{fix}}$ be the following finite-depth event:
\begin{enumerate}[label=(\roman*)]
 \item label $c$ has exactly $m$ points in $Q_c$ and no other point of
       depth at most $R_0$;
 \item label $z$ has exactly $m$ points in $Q_z$, exactly one point in
       $D_{\mathcal A}$, and no other point of depth at most $R_0$;
 \item each label $\ell\in\{x,y\}$ has exactly $m$ points in $Q_\ell$,
       no point in
       $[0,R_0]\setminus(Q_\ell\cup D_{\mathcal K})$, and no condition
       imposed on its process in $D_{\mathcal K}$.
\end{enumerate}
Since $D_{\mathcal A},D_{\mathcal K}\subset(0,1)$ whereas every $Q_\ell$
lies in $(R_0-1,R_0)=(74,75)$, the anchor-depth regions are disjoint from
all four buffer intervals.
We now specify the finite atoms used to evaluate this event.  For label
$c$, take $Q_c$ itself and the nonempty interval components of
$[0,R_0]\setminus Q_c$; prescribe counts $m$ and $0$, respectively.  For
label $z$, take $Q_z$, $D_{\mathcal A}$, and the nonempty interval
components of their complement in $[0,R_0]$; prescribe counts
$m$, $1$, and $0$.  For each $\ell\in\{x,y\}$, take $Q_\ell$ and the
nonempty interval components of
\[
 [0,R_0]\setminus(Q_\ell\cup D_{\mathcal K});
\]
prescribe count $m$ on $Q_\ell$ and count $0$ on every listed component.
The region $D_{\mathcal K}$ is deliberately absent from the constrained
atom family for these two labels.  Endpoint assignments are inherited
from the half-open conventions defining the intervals, so within each
label the listed atoms are disjoint and describe exactly the conditions
in (i)--(iii).  Let $k_{\ell,D}$ denote the prescribed count.  When
$k_{\ell,D}=0$, the corresponding factor below is simply
$\e^{-\Lambda_T(D)}$.
Poisson independent increments give the exact product
\begin{equation}\label{eq:fixed-event-product}
 \widetilde{\mathsf Q}_T(F_{\mathrm{fix}})
 =
 \prod_{(\ell,D)\ \mathrm{constrained}}
 \e^{-\Lambda_T(D)}
 \frac{\Lambda_T(D)^{k_{\ell,D}}}{k_{\ell,D}!}.
\end{equation}
By finite additivity and \eqref{eq:fixed-depth-intensity}, every intensity
in this finite product has a finite limit.  Each atom on which a positive
count is required has positive length and hence a strictly positive
limiting intensity.  It follows that, for all sufficiently large $T$,
\begin{equation}\label{eq:fixed-event-lower}
 \widetilde{\mathsf Q}_T(F_{\mathrm{fix}})
 \ge c_{\mathrm{fix}}>0,
\end{equation}
where $c_{\mathrm{fix}}$ is independent of $T$.

\smallskip
\noindent\emph{Step 3: survival through the entire remaining depth.}
For a label $\ell$, define
\[
 M_{\ell,T}(t):=
 \sum_{\substack{p\in\cB(T)\\R_0<s_p\le R_0+t}}K_{\ell,p},
 \qquad 0\le t\le\lambda_T-R_0.
\]
This independent-increment process can be realized as
\begin{equation}\label{eq:poisson-time-change}
 M_{\ell,T}(t)
 =\Pi_\ell(\Lambda_T((R_0,R_0+t])),
\end{equation}
where the $\Pi_\ell$ are independent rate-one Poisson processes.  Put
$\epsilon_T:=\sup_t|\varepsilon_T(t)|$.  For all sufficiently large $T$,
$q\epsilon_T\le B/2$.

Suppose that
\begin{equation}\label{eq:survival-event}
 \Pi_\ell(v)+B/2\ge qv
 \quad\text{for every }v\ge0.
\end{equation}
With $v=\Lambda_T((R_0,R_0+t])$, the time-change estimate gives
$t/3\le v+\epsilon_T$.  Therefore
\begin{align}
 M_{\ell,T}(t)+B
 &\ge qv+B/2
 \ge \alpha t,                                      \label{eq:profile-growth}\\
 m+M_{\ell,T}(t)
 &\ge \alpha(R_0+t).                                 \label{eq:full-profile-growth}
\end{align}
The second line uses $m\ge\alpha R_0+B$.  By
Lemma~\ref{lem:poisson-survival} and the first inequality in
\eqref{eq:constant-certificate}, the probability that
\eqref{eq:full-profile-growth} fails for at least one of the four labels
is less than $1/8$.

The weighted tails are also small at fixed positive cost.  By
\eqref{eq:tail-mean}, Markov's inequality, and the last inequality in
\eqref{eq:constant-certificate},
\begin{equation}\label{eq:tail-union}
 \Pp\left(
 \max_{\ell\in\{c,x,y,z\}}V_{\ell,T}>\delta/2
 \right)
 \le\frac{8\e^{-R_0}}{\delta}<\frac18.
\end{equation}
All depth-$>R_0$ processes are independent of $F_{\mathrm{fix}}$.  Let
$\mathcal E_T$ be the event that \eqref{eq:full-profile-growth} holds for
all four labels and all $0\le t\le\lambda_T-R_0$, and that
\begin{equation}\label{eq:small-tails}
 V_{\ell,T}\le\delta/2
 \qquad(\ell=c,x,y,z)
\end{equation}
holds.  The preceding two union bounds show that
\begin{equation}\label{eq:profile-tail-success}
 \widetilde{\mathsf Q}_T(\mathcal E_T\mid F_{\mathrm{fix}})
 \ge\frac34.
\end{equation}
No error depending on $w$ has been subtracted.

\smallskip
\noindent\emph{Step 4: two adaptive correcting anchors.}
Let $\mathcal H_T$ be the $\sigma$-field generated by the restrictions of
all four labelled Poisson arrays except for the restrictions of the
$x$- and $y$-arrays to $D_{\mathcal K}$.  The event
$F_{\mathrm{fix}}$ is $\mathcal H_T$-measurable because it imposes no
condition on those two omitted restrictions.  The event $\mathcal E_T$ is
also $\mathcal H_T$-measurable: it is determined entirely by the four
depth-$>R_0$ processes.  By independence across labels and independent
increments on disjoint depth sets, conditionally on $\mathcal H_T$ the
two omitted restrictions remain independent Poisson processes, each with
its original intensity $\Lambda_T$ restricted to $D_{\mathcal K}$.

Let $U_z\in\mathcal A$ be the normalized weight of the unique $z$-point in
$D_{\mathcal A}$.  For $\ell\in\{x,y,z\}$, let $E_\ell$ be the total weight
of the $m$ points in $Q_\ell$ plus $V_{\ell,T}$.  Every point in $Q_\ell$
has weight at most $\e^{1-R_0}$, so
\eqref{eq:constant-certificate} and \eqref{eq:small-tails} imply
\begin{equation}\label{eq:error-weights}
 0\le E_\ell\le\delta.
\end{equation}

Work conditionally on $\mathcal H_T$, on
$F_{\mathrm{fix}}\cap\mathcal E_T$.  The variables
$U_z,E_x,E_y,E_z$ are then fixed.
Define the $\mathcal H_T$-measurable centers
\begin{equation}\label{eq:anchor-centers}
 t_x:=U_z+E_z-E_x,
 \qquad
 t_y:=U_z+E_z-E_y
\end{equation}
and target intervals
\begin{equation}\label{eq:anchor-intervals}
 I_x:=\left[t_x-\frac w8,t_x+\frac w8\right],
 \qquad
 I_y:=\left[t_y-\frac w8,t_y+\frac w8\right].
\end{equation}
From $U_z\in[12/25,13/25)$ and \eqref{eq:error-weights},
\[
 t_x,t_y\in\left[\frac{47}{100},\frac{53}{100}\right].
\]
Since $w\to0$, both intervals in \eqref{eq:anchor-intervals} lie in
$\mathcal K$ for all sufficiently large $T$.

Require the $x$-process to have exactly one point in $I_x$ and no other
point in $D_{\mathcal K}$, and impose the analogous requirement on the
$y$-process.  Let $\mathcal A_T$ denote this two-anchor event on
$F_{\mathrm{fix}}\cap\mathcal E_T$, and declare it empty off that event.
Conditionally on $\mathcal H_T$, the probability of the
$x$-requirement is exactly
\begin{equation}\label{eq:anchor-probability}
 \left(
 \sum_{\substack{p\in\cB(T)\\u_p\in I_x}}h_p
 \right)
 \exp\!\left(
 -\sum_{\substack{p\in\cB(T)\\u_p\in\mathcal K}}h_p
 \right).
\end{equation}
The interval $I_x$ has the form $[t,t+w/4]$ with
$t=t_x-w/8$; for large $T$, this left endpoint lies in the fixed compact
interval $[23/50,27/50]\subset(0,1)$.  Thus
Lemma~\ref{lem:local-prime-mass}, with $c_0=1/4$, gives, almost surely on
$F_{\mathrm{fix}}\cap\mathcal E_T$ and uniformly over every possible
conditional realization,
\[
 \sum_{\substack{p\in\cB(T)\\u_p\in I_x}}h_p
 =\frac13\sum_{\substack{p\in\cB(T)\\u_p\in I_x}}\frac1{p+1}
 \ge c_1w.
\]
Although $I_x$ is random before conditioning, it is
$\mathcal H_T$-measurable and hence is a fixed deterministic interval
under each conditional realization.  The lower bound in
Lemma~\ref{lem:local-prime-mass} is uniform over all left endpoints in
$[23/50,27/50]$, so the same constant $c_1$ applies simultaneously to
every such realization.  The identical observation applies to $I_y$.
Moreover, \eqref{eq:mertens-plus} gives
\[
 \sum_{\substack{p\in\cB(T)\\u_p\in\mathcal K}}h_p
 \longrightarrow\frac13\log\frac{11}{9}.
\]
Consequently the exponential factor in \eqref{eq:anchor-probability} is at
least a fixed $c_2>0$.  The same bounds hold for $I_y$.  Independence of the two
omitted processes therefore gives both correcting anchors with
conditional probability at least
\begin{equation}\label{eq:two-anchor-lower}
 c_3w^2.
\end{equation}

\smallskip
\noindent\emph{Step 5: verification of the target event and transfer.}
Let $U_x,U_y$ be the weights of the two correcting anchors.  On the event
just constructed,
\[
 X=U_x+E_x,\qquad
 Y=U_y+E_y,\qquad
 Z=U_z+E_z.
\]
By \eqref{eq:anchor-centers}--\eqref{eq:anchor-intervals},
\begin{equation}\label{eq:balanced-totals}
 |X-Z|\le w/8,\qquad |Y-Z|\le w/8.
\end{equation}
Thus the range of $X,Y,Z$ is at most $w/4$.  Also
$Z\in[48/100,53/100]$ by \eqref{eq:error-weights}; since $w\to0$,
\eqref{eq:balanced-totals} places all three totals in the fixed closed
interval $J$ for large $T$.

At depth $R_0$, every label has at least $m$ points.  For
$s=R_0+t$, \eqref{eq:full-profile-growth} gives
\[
 N_\ell^{\mathrm P}(s)
 \ge\alpha s\ge\lfloor\alpha s\rfloor
 \qquad(\ell=c,x,y,z).
\]
The constructed event is therefore contained in $B_T^{\mathrm P}$.
The conditional estimate \eqref{eq:two-anchor-lower} is uniform on
$F_{\mathrm{fix}}\cap\mathcal E_T$.  The tower property,
\eqref{eq:fixed-event-lower}, and \eqref{eq:profile-tail-success} therefore
give
\begin{align*}
 \widetilde{\mathsf Q}_T(
 F_{\mathrm{fix}}\cap\mathcal E_T\cap\mathcal A_T)
 &=\E_{\widetilde{\mathsf Q}_T}\!\left[
   \1_{F_{\mathrm{fix}}\cap\mathcal E_T}
   \widetilde{\mathsf Q}_T(\mathcal A_T\mid\mathcal H_T)
   \right]\\
 &\ge c_3w^2\,
   \widetilde{\mathsf Q}_T(F_{\mathrm{fix}}\cap\mathcal E_T)
 \ge\frac34c_{\mathrm{fix}}c_3w^2.
\end{align*}
The event on the left is contained in $B_T^{\mathrm P}$, so
\begin{equation}\label{eq:poisson-first-moment}
 \widetilde{\mathsf Q}_T(B_T^{\mathrm P})
 \ge\frac34c_{\mathrm{fix}}c_3w^2=:c_4w^2.
\end{equation}
Finally, the defining property of total-variation distance gives
\[
 \left|
 \mathsf Q_T(B_T^{\mathrm P})
 -\widetilde{\mathsf Q}_T(B_T^{\mathrm P})
 \right|
 \le d_{\TV}(\mathsf Q_T,\widetilde{\mathsf Q}_T).
\]
Lemma~\ref{lem:poisson-coupling} with $A=2$ makes the right side
$o(w^2)$.  Hence \eqref{eq:event-identification} and
\eqref{eq:poisson-first-moment} give
\[
 \beta_T
 =\mathsf Q_T(B_T^{\mathrm P})
 \ge\widetilde{\mathsf Q}_T(B_T^{\mathrm P})-o(w^2)
 \ge c_4w^2-o(w^2)
 \ge\frac{c_4}{2}w^2
\]
for all sufficiently large $T$.  This proves
\eqref{eq:first-moment}.
\end{proof}

\section{The exact rooted second moment}\label{sec:secondmoment}

The first-moment estimate shows that conditioning on $B_T$ is not too
expensive.  We now show that the conditioned root density does not
concentrate.  The proof is annealed: the root support and its two independent
colourings are summed together.  This is essential, because no comparable
anti-concentration statement is true uniformly for every fixed root.

\subsection{The conditional root law and factorial insertion}

Fix the first state $S=S^{(x)}=c\cup y\cup z$.  Since its marginal is
$\mu_{\cB(T)}$, the primes of $S$ are independent with probabilities
$r_p=1/(p+1)$ and inclusion odds
\begin{equation}\label{eq:root-odds}
 \frac{r_p}{1-r_p}=\frac1p.
\end{equation}
Conditional on $S$, every $p\in S$ is independently and uniformly coloured
$c,y,$ or $z$.  Every $p\notin S$ is independently assigned to the missing
$x$-petal with probability
\begin{equation}\label{eq:missing-probability}
 \frac{h_p}{1-r_p}
 =\frac{1}{3p}.
\end{equation}
Indeed, absence from $S$ means that the original label is either
$\varnothing$ or $x$.

Take two independent completions conditional on the same root.  Put
\[
 \sigma(c)=0,\qquad \sigma(y)=1,\qquad \sigma(z)=-1
\]
and, for $p\in S$, define the nine-valued mark
\begin{equation}\label{eq:nine-marks}
 \xi_p:=(\sigma(\ell_1(p)),\sigma(\ell_2(p)))
 \in\cM:=\{-1,0,1\}^2.
\end{equation}
All nine marks are equiprobable.  If
\begin{equation}\label{eq:root-differences}
 D_h:=Y_h-Z_h\qquad(h=1,2),
\end{equation}
then the vector $D=(D_1,D_2)$ is exactly
\begin{equation}\label{eq:marked-sum}
 D=\sum_{p\in S}u_p\xi_p.
\end{equation}

Both pivot estimates use the following exact identity.

\begin{lemma}[Factorial deletion and insertion]\label{lem:factorial-insertion}
Let $\mathcal P$ be finite, and let $\mu$ be the product law
\[
 \mu(S)=\prod_{p\in S}r_p\prod_{p\notin S}(1-r_p),
 \qquad
 \omega_p:=\frac{r_p}{1-r_p},
 \qquad 0<r_p<1.
\]
For $d\ge1$, let $F(S;p_1,\ldots,p_d)$ be nonnegative and defined when
$p_1,\ldots,p_d$ are distinct members of $S$.  Then
\begin{align}
 &\E_{S\sim\mu}
 \sum_{p_1,\ldots,p_d\in S}^{\ne}
 F(S;p_1,\ldots,p_d)\notag\\
 &\quad=
 \sum_{A\subseteq\mathcal P}\mu(A)
 \sum_{\substack{p_1,\ldots,p_d\in\mathcal P\setminus A\\
                  p_1,\ldots,p_d\ \mathrm{distinct}}}
 \left(\prod_{j=1}^d\omega_{p_j}\right)
 F(A\cup\{p_1,\ldots,p_d\};p_1,\ldots,p_d).
                                                        \label{eq:factorial-insertion}
\end{align}
\end{lemma}

\begin{proof}
Expand the expectation on the left.  For every ordered tuple appearing in
the expansion, make the bijective substitution
\[
 A=S\setminus\{p_1,\ldots,p_d\}.
\]
The deleted points are pairwise distinct and absent from $A$, so
\[
 \mu(A\cup\{p_1,\ldots,p_d\})
 =\mu(A)\prod_{j=1}^d\frac{r_{p_j}}{1-r_{p_j}}.
\]
Substitution gives \eqref{eq:factorial-insertion}.  Every sum is finite, so
no limiting interchange is involved.
\end{proof}

\subsection{Profiles and ordered weights}

Order a nonempty support $S$ by decreasing normalized weight:
\[
 u_{(1)}>u_{(2)}>\cdots>u_{(K)}.
\]
There are no ties because distinct primes have distinct logarithms.

\begin{lemma}[Ordered weights forced by a profile]
\label{lem:ordered-profile}
Suppose that
\begin{equation}\label{eq:root-profile}
 N_S(s):=\#\{p\in S:s_p\le s\}\ge as
 \qquad(R_0\le s\le\lambda_T).
\end{equation}
Then $K=|S|\ge a\lambda_T$, and for $1\le k\le K$,
\begin{equation}\label{eq:ordered-lower}
 u_{(k)}\ge\ell_k,
 \qquad
 \ell_k:=
 \exp\!\left(
 -\max\left\{R_0,\frac{k+1}{a}\right\}
 \right).
\end{equation}
\end{lemma}

\begin{proof}
The inequality $K\ge a\lambda_T$ is
\eqref{eq:root-profile} at $s=\lambda_T$.  Put
$s_{(k)}=-\log u_{(k)}$.  If $s_{(k)}\le R_0$, then
$u_{(k)}\ge\e^{-R_0}\ge\ell_k$.  Otherwise take
$R_0\le s<s_{(k)}$ and let $s\uparrow s_{(k)}$.  At depth $s$ there are at
most $k-1$ points, so $as\le k-1$.  Hence
$s_{(k)}\le(k-1)/a$, which is stronger than
\eqref{eq:ordered-lower}.
\end{proof}

We record the two convergent endpoint estimates before using them.  From
\eqref{eq:parameter-inequalities},
\begin{equation}\label{eq:endpoint-decay}
 \frac{3^{-a\lambda_T}}{w}
 =
 \frac{T^{\,1-a(1-\vartheta)\log3}}{\eta}
 =o(1).
\end{equation}
Indeed the exponent of $T$ is a fixed negative number and
$1/\eta=T^{o(1)}$.  Moreover,
\begin{equation}\label{eq:rank-series}
 \sum_{k\ge1}\frac{3^{-k}}{\ell_k}<\infty.
\end{equation}
To see this, discard finitely many indices controlled by $R_0$; the
remaining ratio is $\e^{1/a}/3<1$, which is equivalent to
$a\log3>1$.

\begin{lemma}[Annealed two-pivot small-ball bound]
\label{lem:two-pivot}
Let $S\sim\mu_{\cB(T)}$, give its primes independent uniform marks in
$\cM=\{-1,0,1\}^2$, and define $D$ by \eqref{eq:marked-sum}.  For every
axis-parallel square $Q\subseteq\mathbb R^2$ of side at most $C_Qw$,
uniformly in its position,
\begin{equation}\label{eq:two-pivot}
 \Pp\bigl(
 N_S(s)\ge as\text{ for }R_0\le s\le\lambda_T,
 \ D\in Q
 \bigr)
 \ll_{C_Q}w^2.
\end{equation}
\end{lemma}

\begin{proof}
Write
\[
 \mathcal P_T(S):=
 \{N_S(s)\ge as\text{ for every }R_0\le s\le\lambda_T\}.
\]
For brevity, write $\mu:=\mu_{\cB(T)}$.
For each realization put $K:=|S|$.
All ranks below refer to decreasing order of the weights $u_p$.  For a
marked realization $\xi=(\xi_p)_{p\in S}$, define
\[
 r(\xi):=\dim_{\mathbb R}\operatorname{span}\{\xi_p:p\in S\}.
\]
The span of the empty set is understood to be $\{0\}$.  Thus the cases
$r(\xi)=0,1,2$ are disjoint and exhaustive.

\smallskip
\noindent\emph{Mark-counting sublemma.}
For $v\in\cM\setminus\{0\}$, put
\[
 L(v):=\cM\cap\mathbb Rv=\{-v,0,v\}.
\]
Thus $|L(v)|=3$ and $|\cM\setminus L(v)|=6$.  On a support of size $K$,
the total probability of all mark sequences whose first nonzero mark has
rank $i$ and whose first subsequent mark outside its real span has rank
$j>i$ is
\begin{equation}\label{eq:rank-two-mark-mass}
 9^{-K}\cdot8\cdot3^{j-i-1}\cdot6\cdot9^{K-j}
 =16\,3^{-(i+j)}.
\end{equation}
The total probability of all rank-one mark sequences whose first nonzero
mark has rank $i$ is
\begin{equation}\label{eq:rank-one-mark-mass}
 9^{-K}\cdot8\cdot3^{K-i}
 =8\,3^{-(K+i)}.
\end{equation}
Indeed, there are eight choices for the mark at rank $i$.  For
\eqref{eq:rank-two-mark-mass}, the ranks $i+1,\ldots,j-1$ have the three
choices in $L(v)$, rank $j$ has the six choices outside $L(v)$, and every
later rank has nine choices.  For \eqref{eq:rank-one-mark-mass}, every
later mark must belong to $L(v)$.  Multiplication and division by the
$9^K$ equally likely mark sequences prove both formulas.

\smallskip
\noindent\emph{Rank two.}
For a rank-two marked support, let $i$ be the rank of its first nonzero
mark and let $j>i$ be the rank of its first mark outside the span of the
mark at rank $i$.  Let $p$ and $q$ be the primes at these ranks.  This
choice is canonical.

We spell out the data remaining after deletion of $p$ and $q$.  Put
$A=S\setminus\{p,q\}$ and $K=|A|+2$.  Order $A$ as
$r_1,\ldots,r_{K-2}$ with
\[
 b_1:=u_{r_1}>\cdots>b_{K-2}:=u_{r_{K-2}},
 \qquad b_0:=+\infty,\quad b_{K-1}:=0.
\]
When $A=\varnothing$, the displayed ordered list is empty and only the
two endpoint conventions apply.
For $1\le i<j\le K$, let $\mathcal G_{ij}(A)$ be the set of ordered pairs
of distinct primes $p,q\in\cB(T)\setminus A$ such that
\begin{equation}\label{eq:rank-two-gaps}
 \begin{cases}
 b_{i-1}>u_p>u_q>b_i, & j=i+1,\\
 b_{i-1}>u_p>b_i\ \text{ and }\
 b_{j-2}>u_q>b_{j-1}, & j\ge i+2.
 \end{cases}
\end{equation}
These inequalities say exactly that $p$ and $q$ occupy ranks $i$ and $j$
after insertion.  For $v\ne0$, let $\mathcal C_{ij}(A;v)$ be the set of
mark vectors $\zeta=(\zeta_1,\ldots,\zeta_{K-2})$ on this ordered copy of
$A$ satisfying
\begin{equation}\label{eq:rank-two-residual-marks}
 \zeta_k=0\ (k<i),\qquad
 \zeta_k\in L(v)\ (i\le k\le j-2),\qquad
 \zeta_k\in\cM\ (k\ge j-1),
\end{equation}
where an empty index range imposes no condition.  Put
\[
 D_A(\zeta):=\sum_{k=1}^{K-2}b_k\zeta_k.
\]

We now verify explicitly that deletion and insertion preserve every
canonical index.  Let $t_1,\ldots,t_K$ denote the decreasing-weight
ordering after $p,q$ are inserted.  For
$(p,q)\in\mathcal G_{ij}(A)$ this ordering is, and can only be,
\begin{equation}\label{eq:rank-two-reconstruction}
 (t_1,\ldots,t_K)=
 (r_1,\ldots,r_{i-1},p,r_i,\ldots,r_{j-2},q,
  r_{j-1},\ldots,r_{K-2}),
\end{equation}
where every sublist with lower index larger than its upper index is empty.
If $j\ge i+2$, the two insertion conditions in
\eqref{eq:rank-two-reconstruction} are precisely
\[
 b_{i-1}>u_p>b_i,\qquad b_{j-2}>u_q>b_{j-1}.
\]
If $j=i+1$, both points lie in the same residual gap and their required
order is precisely
\[
 b_{i-1}>u_p>u_q>b_i.
\]
Thus \eqref{eq:rank-two-reconstruction} is equivalent to
$(p,q)\in\mathcal G_{ij}(A)$, including the endpoint ranks under the
sentinel conventions.

For an ordered marked support $(S,\xi)$, define
$\mathcal E_{ij}^{v,v'}(S,\xi;p,q)$ to be the event
\[
 \begin{gathered}
 p=t_i,\qquad q=t_j,\qquad
 \xi_{t_k}=0\quad(k<i),\\
 \xi_p=v\ne0,\qquad
 \xi_{t_k}\in L(v)\quad(i<k<j),\qquad
 \xi_q=v'\notin L(v).
 \end{gathered}
\]
There is no restriction on marks after rank $j$.  Pointwise in
$(S,\xi)$,
\begin{equation}\label{eq:rank-two-canonical-identity}
 \1_{\{r(\xi)=2\}}
 =
 \sum_{\substack{p,q\in S\\p\ne q}}
 \sum_{1\le i<j\le K}
 \sum_{v\in\cM\setminus\{0\}}
 \sum_{v'\in\cM\setminus L(v)}
 \1_{\mathcal E_{ij}^{v,v'}(S,\xi;p,q)}.
\end{equation}
Indeed, if $r(\xi)=2$, the sole nonzero summand uses the prime carrying
the first nonzero mark and the prime carrying the first later mark outside
its span.  If $r(\xi)<2$, no summand can occur.  In particular, exchanging
$p$ and $q$ never creates a second summand: their canonical ranks are
different, and \eqref{eq:rank-two-reconstruction} assigns each rank to one
specific insertion gap.

Under \eqref{eq:rank-two-reconstruction}, deletion of the two pivot marks
turns the full-rank conditions in
$\mathcal E_{ij}^{v,v'}$ into exactly
\eqref{eq:rank-two-residual-marks}.  Namely, residual indices $k<i$
retain full rank $k$; indices $i\le k\le j-2$ retain full rank $k+1$ and
lie strictly between the pivots; and indices $k\ge j-1$ retain full rank
$k+2$ and occur after the second pivot.

To apply Lemma~\ref{lem:factorial-insertion} without suppressing the mark
average, for distinct $p,q\in S$ define
\begin{align*}
 F_2(S;p,q)
 &:=
 9^{-|S|}\sum_{\xi\in\cM^S}
 \1_{\mathcal P_T(S)}
 \1_{\{\sum_{s\in S}u_s\xi_s\in Q\}}\\
 &\qquad\times
 \sum_{1\le i<j\le |S|}
 \sum_{v\in\cM\setminus\{0\}}
 \sum_{v'\in\cM\setminus L(v)}
 \1_{\mathcal E_{ij}^{v,v'}(S,\xi;p,q)}.
\end{align*}
Equation~\eqref{eq:rank-two-canonical-identity} gives the exact identity
\[
 \Pp(r(\xi)=2,\mathcal P_T(S),D\in Q)
 =
 \E_{S\sim\mu}\sum_{\substack{p,q\in S\\p\ne q}}F_2(S;p,q).
\]
For $p,q\notin A$, the support and mark coefficient after deletion is
\begin{align*}
 \mu(A\cup\{p,q\})9^{-(|A|+2)}
 &=\mu(A)
 \frac{1/(p+1)}{1-1/(p+1)}
 \frac{1/(q+1)}{1-1/(q+1)}
 9^{-(|A|+2)}\\
 &=\frac{\mu(A)}{pq}\,9^{-(|A|+2)}.
\end{align*}
Thus the two-point case of
Lemma~\ref{lem:factorial-insertion}, followed by the unique decomposition
of the marks into $(v,v',\zeta)$ described above, gives the exact equality
\begin{align}
 &\Pp(r(\xi)=2,\ \mathcal P_T(S),\ D\in Q)\notag\\
 &=\sum_{A\subseteq\cB(T)}\mu(A)
   \sum_{1\le i<j\le |A|+2}9^{-(|A|+2)}
   \sum_{v\in\cM\setminus\{0\}}
   \sum_{v'\in\cM\setminus L(v)}
   \sum_{\zeta\in\mathcal C_{ij}(A;v)}\notag\\
 &\quad\times
   \sum_{(p,q)\in\mathcal G_{ij}(A)}\frac1{pq}
   \1_{\mathcal P_T(A\cup\{p,q\})}
   \1_{\{D_A(\zeta)+u_pv+u_qv'\in Q\}}.
                                                        \label{eq:rank-two-exact}
\end{align}
Here factorial insertion changes the support mass from
$\mu(A\cup\{p,q\})$ to $\mu(A)/(pq)$;
\eqref{eq:rank-two-gaps} restores the two deleted ranks; and
\eqref{eq:rank-two-residual-marks} restores exactly the residual mark
restrictions.  Conversely, every tuple counted on the right reconstructs
a rank-two marked support whose canonical data are
$(i,j,p,q,v,v')$.  Thus \eqref{eq:rank-two-exact} has neither omission nor
multiplicity.

Fix a term through the choice of $A,i,j,v,v',\zeta$.  On the profile
event, Lemma~\ref{lem:ordered-profile} gives
\begin{equation}\label{eq:pivot-lower-weights}
 u_p\ge\ell_i,\qquad u_q\ge\ell_j.
\end{equation}
Let $M_{v,v'}$ be the matrix with columns $v,v'$.  It has nonzero integer
determinant.  Every cofactor has absolute value at most one, while the
determinant has absolute value at least one; hence every entry of
$M_{v,v'}^{-1}$ has absolute value at most one.  Define
\[
 \mathfrak R:=M_{v,v'}^{-1}(Q-D_A(\zeta)),
 \qquad J_h:=\pi_h(\mathfrak R)\quad(h=1,2).
\]
Each $J_h$ is an interval of length at most $2C_Qw$.  Indeed, if
$z,z'\in Q$ and $\rho_h$ is row $h$ of $M_{v,v'}^{-1}$, then
\[
 |\rho_h(z-z')|
 \le\|\rho_h\|_1\max_{k=1,2}|z_k-z'_k|
 \le2C_Qw.
\]

Let $H_i,H_j$ denote the two individual rank gaps in
\eqref{eq:rank-two-gaps}; when $j=i+1$, both equal
$(b_i,b_{i-1})$, while for $j\ge i+2$ they are
\[
 H_i=(b_i,b_{i-1}),\qquad H_j=(b_{j-1},b_{j-2}).
\]
Taking closures only enlarges the sets, so define
\[
 I_i:=\overline{
 J_1\cap H_i\cap[\ell_i,1]\cap[T^{\vartheta-1},1]},
 \qquad
 I_j:=\overline{
 J_2\cap H_j\cap[\ell_j,1]\cap[T^{\vartheta-1},1]}.
\]
If either intersection is empty, the corresponding term in
\eqref{eq:rank-two-exact} is zero.  Otherwise
$I_i,I_j\subseteq[T^{\vartheta-1},1]$, both have length at most $2C_Qw$,
and
\[
 \inf I_i\ge\ell_i,\qquad \inf I_j\ge\ell_j.
\]
Every pair retained by \eqref{eq:rank-two-exact} lies in the Cartesian
product determined by $I_i,I_j$.  Hence
\begin{align}
 &\sum_{(p,q)\in\mathcal G_{ij}(A)}\frac1{pq}
   \1_{\mathcal P_T(A\cup\{p,q\})}
   \1_{\{D_A(\zeta)+u_pv+u_qv'\in Q\}}\notag\\
 &\quad\le
 \left(\sum_{\substack{p\in\cB(T)\\u_p\in I_i}}\frac1p\right)
 \left(\sum_{\substack{q\in\cB(T)\\u_q\in I_j}}\frac1q\right)
 \ll_{C_Q}\frac{w^2}{\ell_i\ell_j}.
                                                        \label{eq:two-local-sums}
\end{align}
The product on the right deliberately adds pairs with $p=q$, primes
belonging to $A$, the wrong order $u_p\le u_q$ in the adjacent-rank case,
and points of $J_1\times J_2$ outside the parallelogram $\mathfrak R$.  It
also discards the profile restriction.  All summands are nonnegative, so
each enlargement has the required upper-bound direction.

The final bound in \eqref{eq:two-local-sums} is uniform in
$A,i,j,v,v',\zeta$.  The mark-counting sublemma therefore allows the mark
sums in \eqref{eq:rank-two-exact} to be performed before the remaining
restrictions are discarded:
\begin{align}
 &\Pp(r(\xi)=2,\ \mathcal P_T(S),\ D\in Q)\notag\\
 &\quad\ll_{C_Q}w^2
 \sum_{A\subseteq\cB(T)}\mu(A)
 \sum_{1\le i<j\le |A|+2}
 \frac{16\,3^{-(i+j)}}{\ell_i\ell_j}\notag\\
 &\quad\le
 16Cw^2
 \left(\sum_{k\ge1}\frac{3^{-k}}{\ell_k}\right)^2
 \sum_{A\subseteq\cB(T)}\mu(A)
 \ll_{C_Q}w^2,                                      \label{eq:rank-two-bound}
\end{align}
where \eqref{eq:rank-series} and $\sum_A\mu(A)=1$ were used.

\smallskip
\noindent\emph{Rank one.}
Let $p$ be the canonical first nonzero pivot and let $i$ be its rank.
After deleting $p$, put $A=S\setminus\{p\}$ and $K=|A|+1$.  Order $A$
with weights $b_1>\cdots>b_{K-1}$ and put
$b_0=+\infty$, $b_K=0$.  The pivot has rank $i$ exactly when
\[
 u_p\in G_i(A):=(b_i,b_{i-1}).
\]
For $v\ne0$, let $\mathcal C_i^{(1)}(A;v)$ consist of the residual marks
$\zeta=(\zeta_1,\ldots,\zeta_{K-1})$ satisfying
\[
 \zeta_k=0\quad(k<i),\qquad
 \zeta_k\in L(v)\quad(k\ge i).
\]
When $A=\varnothing$, this residual list is empty and the endpoint
conventions still apply.  Put
\[
 D_A(\zeta):=\sum_{k=1}^{K-1}b_k\zeta_k.
\]
Let $t_1,\ldots,t_K$ denote the reconstructed decreasing-weight order.
It is uniquely
\begin{equation}\label{eq:rank-one-reconstruction}
 (t_1,\ldots,t_K)
 =(r_1,\ldots,r_{i-1},p,r_i,\ldots,r_{K-1}),
\end{equation}
which is equivalent to $b_{i-1}>u_p>b_i$.  For an ordered marked support
$(S,\xi)$, let $\mathcal E_i^v(S,\xi;p)$ denote the event
\[
 p=t_i,\qquad
 \xi_{t_k}=0\ (k<i),\qquad
 \xi_p=v\ne0,\qquad
 \xi_{t_k}\in L(v)\ (k>i).
\]
Then, pointwise,
\begin{equation}\label{eq:rank-one-canonical-identity}
 \1_{\{r(\xi)=1\}}
 =
 \sum_{p\in S}\sum_{i=1}^K
 \sum_{v\in\cM\setminus\{0\}}
 \1_{\mathcal E_i^v(S,\xi;p)}.
\end{equation}
The unique nonzero summand uses the first nonzero marked prime.  Deleting
that prime from \eqref{eq:rank-one-reconstruction} makes residual indices
$k<i$ the marks before the pivot and indices $k\ge i$ the marks after it.
Thus the residual conditions are exactly those defining
$\mathcal C_i^{(1)}(A;v)$, and in particular
$|\mathcal C_i^{(1)}(A;v)|=3^{K-i}$.

For $p\in S$, define
\begin{align*}
 F_1(S;p)
 &:=
 9^{-|S|}\sum_{\xi\in\cM^S}
 \1_{\mathcal P_T(S)}
 \1_{\{\sum_{s\in S}u_s\xi_s\in Q\}}
 \sum_{i=1}^{|S|}
 \sum_{v\in\cM\setminus\{0\}}
 \1_{\mathcal E_i^v(S,\xi;p)}.
\end{align*}
Equation~\eqref{eq:rank-one-canonical-identity} gives
\[
 \Pp(r(\xi)=1,\mathcal P_T(S),D\in Q)
 =\E_{S\sim\mu}\sum_{p\in S}F_1(S;p).
\]
For $p\notin A$, the coefficient after deletion is exactly
\[
 \mu(A\cup\{p\})9^{-(|A|+1)}
 =\mu(A)\frac{1/(p+1)}{1-1/(p+1)}9^{-(|A|+1)}
 =\frac{\mu(A)}p\,9^{-(|A|+1)}.
\]
The one-point case of Lemma~\ref{lem:factorial-insertion}, followed by the
unique decomposition of the remaining marks as $\zeta$, therefore gives
the exact equality
\begin{align}
 &\Pp(r(\xi)=1,\ \mathcal P_T(S),\ D\in Q)\notag\\
 &=\sum_{A\subseteq\cB(T)}\mu(A)
   \sum_{i=1}^{|A|+1}9^{-(|A|+1)}
   \sum_{v\in\cM\setminus\{0\}}
   \sum_{\zeta\in\mathcal C_i^{(1)}(A;v)}
   \sum_{\substack{p\in\cB(T)\setminus A\\u_p\in G_i(A)}}\frac1p
   \1_{\mathcal P_T(A\cup\{p\})}
   \1_{\{D_A(\zeta)+u_pv\in Q\}}.
                                                        \label{eq:rank-one-exact}
\end{align}

Fix $A,i,v,\zeta$ and choose a coordinate $h$ with $|v_h|=1$.  Write
$Q=Q_1\times Q_2$.  The last indicator in
\eqref{eq:rank-one-exact} forces
\[
 u_p\in J_i^0:=v_h(Q_h-D_{A,h}(\zeta)),
\]
an interval of length at most $C_Qw$.  On the profile event,
Lemma~\ref{lem:ordered-profile} also gives $u_p\ge\ell_i$.  Intersect
$J_i^0$ with the rank gap, the profile lower bound, and the prime band,
and set
\[
 I_i^{(1)}:=
 \overline{
 J_i^0\cap G_i(A)\cap[\ell_i,1]\cap[T^{\vartheta-1},1]}.
\]
If the intersection is empty, the term is zero.  Otherwise
$I_i^{(1)}\subseteq[T^{\vartheta-1},1]$ is an interval of length at most
$C_Qw$ and $\inf I_i^{(1)}\ge\ell_i$.  Therefore
\begin{align}
 &\sum_{\substack{p\in\cB(T)\setminus A\\u_p\in G_i(A)}}\frac1p
   \1_{\mathcal P_T(A\cup\{p\})}
   \1_{\{D_A(\zeta)+u_pv\in Q\}}\notag\\
 &\qquad\le
 \sum_{\substack{p\in\cB(T)\\u_p\in I_i^{(1)}}}\frac1p
 \ll_{C_Q}\frac{w}{\ell_i}.                         \label{eq:one-local-sum}
\end{align}
The right side adds back primes in $A$ and discards the profile and the
unused coordinate condition from $D_A(\zeta)+u_pv\in Q$; nonnegativity
gives the displayed upper-bound direction.

The profile event further implies $K\ge a\lambda_T$.  We may now drop all
of its other restrictions and sum the residual marks explicitly.  The
mark-counting sublemma gives
\begin{align}
 &\Pp(r(\xi)=1,\ \mathcal P_T(S),\ D\in Q)\notag\\
 &\quad\le Cw\sum_{A\subseteq\cB(T)}\mu(A)
   \1_{\{|A|+1\ge a\lambda_T\}}
   \sum_{i=1}^{|A|+1}
   \frac{8\,3^{-(|A|+1+i)}}{\ell_i}\notag\\
 &\quad\le
 8Cw\,3^{-a\lambda_T}
 \left(\sum_{i\ge1}\frac{3^{-i}}{\ell_i}\right)
 \sum_{A\subseteq\cB(T)}\mu(A)
 =o(w^2),                                           \label{eq:rank-one-bound}
\end{align}
by \eqref{eq:rank-series}, \eqref{eq:endpoint-decay}, and
$\sum_A\mu(A)=1$.

\smallskip
\noindent\emph{Rank zero.}
All $K$ marks vanish, which has conditional probability $9^{-K}$.
The profile gives $K\ge a\lambda_T$, so
\begin{equation}\label{eq:rank-zero-bound}
 \Pp(r(\xi)=0,\ \mathcal P_T(S),\ D\in Q)
 \le9^{-a\lambda_T}
 =o(w^2)
\end{equation}
by \eqref{eq:endpoint-decay}.  The three rank cases prove
\eqref{eq:two-pivot}.
\end{proof}

The finite counts in \eqref{eq:rank-two-mark-mass} and
\eqref{eq:rank-one-mark-mass}, as well as the inverse-matrix bound for all
48 ordered noncollinear mark pairs, are independently enumerated by the
companion verifier.

\subsection{The missing-petal pivot}

\begin{lemma}[Canonical top missing-petal pivot]
\label{lem:missing-pivot}
Fix a root support $S\subseteq\cB(T)$.  Independently select each
$p\in\cB(T)\setminus S$ with probability $q_p=1/(3p)$, and let
\[
 X=\sum_{\text{selected }p}u_p.
\]
For every interval $I$ of length at most $C_Xw$,
\begin{equation}\label{eq:missing-one-copy}
 \Pp\bigl(X\in I,\ N_x(R_0)\ge1\mid S\bigr)
 \ll_{C_X}w,
\end{equation}
uniformly in $S$ and $I$.  For two conditionally independent completions
and two target intervals fixed after the root data have been exposed, the
corresponding joint probability is $O_{C_X}(w^2)$.
\end{lemma}

\begin{proof}
On the event in \eqref{eq:missing-one-copy}, choose canonically the selected
prime $p_*$ of largest weight among those of depth at most $R_0$, and delete
it.  Let $\nu$ be the product law with parameters $q_p$ on
$\cB(T)\setminus S$.  The one-point case of
Lemma~\ref{lem:factorial-insertion} gives the exact identity
\begin{align}
 &\Pp\bigl(X\in I,\ N_x(R_0)\ge1\mid S\bigr)\notag\\
 &\quad=
 \sum_{A\subseteq\cB(T)\setminus S}\nu(A)
 \sum_{p\in(\cB(T)\setminus S)\setminus A}
 \frac{q_p}{1-q_p}\,
 \1_{\{p\ \mathrm{is\ canonical}\}}
 \1_{\{u_p+\sum_{r\in A}u_r\in I\}}.                  \label{eq:missing-insertion}
\end{align}
There is exactly one canonical point on the event, so no multiplicity is
present in this identity.

The canonical restriction gives $u_p\ge\e^{-R_0}$.  For fixed $A$, the last
indicator in \eqref{eq:missing-insertion} confines $u_p$ to a translate of
$I$.  Moreover,
\begin{equation}\label{eq:missing-odds}
 \frac{q_p}{1-q_p}=\frac1{3p-1}\le\frac1p.
\end{equation}
Discard the remaining canonical order restriction and apply
Lemma~\ref{lem:local-prime-mass} on $[\e^{-R_0},1]$.  The inner sum in
\eqref{eq:missing-insertion} is $O_{C_X}(w)$ uniformly in $A,S$.
Summing $\nu(A)$ proves \eqref{eq:missing-one-copy}.

Given the root data, the two missing-$x$ processes are independent, so the
two one-copy estimates multiply.  The product already includes the
possibility that both copies select the same prime; no diagonal correction
is needed.
\end{proof}

\subsection{The second moment and rooted balance}

\begin{corollary}[Rooted second moment]\label{cor:rooted-second-moment}
For the function $g_T$ in \eqref{eq:rooted-functions},
\begin{equation}\label{eq:second-moment}
 \E_{S\sim\mu_{\cB(T)}}g_T(S)^2
 =\Pp(B_T^{(1)}\cap B_T^{(2)})
 \ll w^4.
\end{equation}
\end{corollary}

\begin{proof}
Let $B_T^{(1)}$ and $B_T^{(2)}$ be formed from two independent
completions conditional on the same root $S$.  Then, conditionally on $S$,
\[
 \Pp(B_T^{(1)}\cap B_T^{(2)}\mid S)
 =\Pp(B_T\mid S)^2=g_T(S)^2.
\]
Taking expectations proves the equality in \eqref{eq:second-moment}.

Let
\[
 \mathcal F_{\mathrm{root}}
 :=\sigma\bigl(S,(\xi_p)_{p\in S}\bigr),
\]
where the two coordinates of $\xi_p$ record the two complete colourings
of the root support as in \eqref{eq:nine-marks}.  Thus
$\mathcal F_{\mathrm{root}}$ reveals the common root and both root
colourings, but neither missing-$x$ process.  Define the root-good event
\[
 \mathcal R_T:=
 \left\{
 N_S(s)\ge as\text{ for every }R_0\le s\le\lambda_T
 \right\}
 \cap\{D\in[-w,w]^2\}.
\]
It is $\mathcal F_{\mathrm{root}}$-measurable.  If both completions satisfy
$B_T$, then, in completion $h$, the root labels $c,y,z$ partition $S$.
Consequently, for $R_0\le s\le\lambda_T$,
\[
 N_S(s)
 =N_{c,h}(s)+N_{y,h}(s)+N_{z,h}(s)
 \ge3\lfloor\alpha s\rfloor
 \ge as
\]
by \eqref{eq:profile-dominates}.  Moreover
$D_h=Y_h-Z_h\in[-w,w]$.  Hence
$B_T^{(1)}\cap B_T^{(2)}\subseteq\mathcal R_T$, and
Lemma~\ref{lem:two-pivot}, with $Q=[-w,w]^2$, gives
\begin{equation}\label{eq:root-good-probability}
 \Pp(\mathcal R_T)\le C_0w^2
\end{equation}
for an absolute constant $C_0$.

For $h=1,2$, let $\mathcal X_h\subseteq\cB(T)\setminus S$ be the primes
assigned to the missing $x$-petal in completion $h$.  By
\eqref{eq:missing-probability}, for every fixed realization of
$\mathcal F_{\mathrm{root}}$ the indicators
\[
 \left(
 \1_{\{p\in\mathcal X_h\}}:
 p\in\cB(T)\setminus S,\ h\in\{1,2\}
 \right)
\]
are mutually independent and satisfy
\[
 \Pp(p\in\mathcal X_h\mid\mathcal F_{\mathrm{root}})=\frac1{3p}.
\]
Indeed, conditional on $S$, the two completions are independent, and
within each completion the colouring variables on $S$ and the
missing-$x$ variables on $\cB(T)\setminus S$ are disjoint independent
factors.  Thus further conditioning on the two root colourings leaves the
two missing processes independent and leaves each with exactly the
product law used in Lemma~\ref{lem:missing-pivot}.

The variables $Y_h,Z_h$ are $\mathcal F_{\mathrm{root}}$-measurable.  On
$\mathcal R_T$, define the corresponding
$\mathcal F_{\mathrm{root}}$-measurable interval
\begin{equation}\label{eq:missing-target}
 I_h:=[\max(Y_h,Z_h)-w,\ \min(Y_h,Z_h)+w],
\end{equation}
and define $I_h:=\{0\}$ off $\mathcal R_T$.  On $\mathcal R_T$ one has
$|Y_h-Z_h|\le w$, so $I_h$ is nonempty and
\[
 |I_h|=2w-|Y_h-Z_h|\le2w.
\]
The range condition in $B_T^{(h)}$ forces $X_h\in I_h$, whereas its
$x$-profile at $R_0$ forces
\[
 N_{x,h}(R_0)\ge\lfloor\alpha R_0\rfloor=24\ge1.
\]
Put
\[
 E_h:=\{X_h\in I_h,\ N_{x,h}(R_0)\ge1\}.
\]
We have the explicit event inclusion
\begin{equation}\label{eq:second-moment-event-inclusion}
 B_T^{(1)}\cap B_T^{(2)}
 \subseteq\mathcal R_T\cap E_1\cap E_2.
\end{equation}
All the other window and profile requirements in
$B_T^{(1)}\cap B_T^{(2)}$ have been discarded on the right; because an
upper bound is sought, this only enlarges the event.

After conditioning on $\mathcal F_{\mathrm{root}}$, the root $S$ and the
two intervals $I_h$ are fixed.  On $\mathcal R_T$ their lengths are at
most $2w$.  The uniform one-copy estimate in
Lemma~\ref{lem:missing-pivot} and the conditional independence of
$\mathcal X_1,\mathcal X_2$ therefore give
\[
 \Pp(E_1\cap E_2\mid\mathcal F_{\mathrm{root}})
 =\prod_{h=1}^2
   \Pp(E_h\mid\mathcal F_{\mathrm{root}})
 \le C_1^2w^2
\]
on $\mathcal R_T$.  Finally,
\eqref{eq:second-moment-event-inclusion}, the tower property, and
\eqref{eq:root-good-probability} yield
\begin{align*}
 \Pp(B_T^{(1)}\cap B_T^{(2)})
 &\le
 \E\!\left[
  \1_{\mathcal R_T}
  \Pp(E_1\cap E_2\mid\mathcal F_{\mathrm{root}})
 \right]\\
 &\le C_1^2w^2\Pp(\mathcal R_T)
 \le C_0C_1^2w^4.
\end{align*}
This proves the upper bound in \eqref{eq:second-moment}.
\end{proof}

\begin{proposition}[Uniform rooted category balance]
\label{prop:rooted-balance}
There is a constant $K<\infty$ with the following uniformity property.
For every function $\eta(T)$ satisfying \eqref{eq:slow-eta}, and all
sufficiently large $T$,
\begin{equation}\label{eq:rooted-L2}
 \E_{S\sim\mu_{\cB(T)}}
 \left(\frac{g_T(S)}{\beta_T}\right)^2
 \le K.
\end{equation}
The constant $K$ is independent of $T$, of $\eta$, of the choice of a band
at a later stage, and of the later integers $H$ and $M$.
\end{proposition}

\begin{proof}
Corollary~\ref{cor:rooted-second-moment} gives a constant $C_2$ such that
the numerator in
\[
 \E\left(\frac{g_T}{\beta_T}\right)^2
 =\frac{\E g_T^2}{\beta_T^2}
\]
is at most $C_2w^4$.  Proposition~\ref{prop:first-moment} gives a constant
$c_1>0$, independent of $T$ and $\eta$, such that
$\beta_T\ge c_1w^2$.  Therefore
\[
 \E\left(\frac{g_T}{\beta_T}\right)^2
 \le\frac{C_2}{c_1^2}.
\]
Take $K=C_2/c_1^2$.  All structural and buffer constants were fixed before
$T$ and $\eta$ varied, so this $K$ has the asserted uniformity.
\end{proof}

\section{Flattening by alternative bands}\label{sec:flattening}

The conditioned law on one band has a uniformly bounded $L^2$ density, but
Proposition~\ref{prop:joint-prefix} needs every word marginal to be close to
the ambient product law in $L^1$.  We obtain that stronger conclusion by
letting each cube coordinate choose its active band uniformly among many
independent alternatives.

\begin{proposition}[Alternative-band flattening]
\label{prop:alternative-bands}
Fix positive integers $H,M$.  For each $1\le i\le H$, let
\[
 \cB_{i,1},\ldots,\cB_{i,M}
\]
be mutually disjoint bands of the form
$\cB_{i,j}=\cB(T_{i,j})$.  Choose a balancing parameter
$\eta_{i,j}>0$ for each band, put
\begin{equation}\label{eq:alternative-band-data}
 w_{i,j}:=\frac{\eta_{i,j}}{T_{i,j}},
\end{equation}
and let $B_{i,j}$ denote the event \eqref{eq:target-event} with
$(T,\eta,w)$ replaced by
$(T_{i,j},\eta_{i,j},w_{i,j})$.  Assume that every $T_{i,j}$ is in the
valid range of Sections~\ref{sec:firstmoment}--\ref{sec:secondmoment}, that
the same constant $K$ in \eqref{eq:rooted-L2} works on all these bands, and
that $\eta_{i,j}\le\eta_*$ for all $i,j$.  Let $E$ be any finite prime set
containing all $HM$ bands.

Then there is a probability law on dimension-$H$ pair-product cubes in
$2^E$ such that:
\begin{enumerate}[label=(\roman*)]
 \item every petal is nonempty;
 \item for every word $\omega$, if $G_\omega$ is the density of its
       marginal relative to $\mu_E$, then
       \begin{equation}\label{eq:flattened-L1}
        \E_{\mu_E}|G_\omega-1|
        \le H\sqrt{\frac KM};
       \end{equation}
 \item every sampled cube satisfies
       \begin{equation}\label{eq:flattened-balance}
        \left|
        \log d(S_\omega)
        -\frac1{3^H}\sum_{\tau\in\F_3^H}\log d(S_\tau)
        \right|
        \le H\eta_*
       \end{equation}
       for every word $\omega$.
\end{enumerate}
\end{proposition}

\begin{proof}
For every group $i$, choose independently an active index
$J_i$ uniformly from $\{1,\ldots,M\}$.  On the active band
$\cB_{i,J_i}$, sample the five-state law conditioned on its event
$B_{i,J_i}$.
Put its $c$-labelled primes into the common support and use its
$x,y,z$ labels as the three petals of coordinate $i$.  On every inactive
band, select each prime with probability $1/(p+1)$ and put every selected
prime into the common support.  Do the same for primes of $E$ outside all
candidate bands.  All choices on disjoint bands are independent.

The bands are disjoint, so petals from different coordinates are disjoint.
Within an active band, the five-state law gives mutually disjoint petals.
Finally, the profile at $R_0$ supplies at least
$\lfloor\alpha R_0\rfloor=24$ points in each petal.  Thus the construction
is a genuine pair-product cube and proves (i).

We next compute a word marginal.  Fix one group $i$ and one of the three
coordinate states selected by the word.  For candidate band $j$, write
$g_{i,j}(S):=\Pp(B_{i,j}\mid S^{(x)}=S)$ and
$\beta_{i,j}:=\Pp(B_{i,j})$, as in
\eqref{eq:rooted-functions}.

We verify here that this density is the same for all three states.  For
$r\in\{x,y,z\}$, put
\[
 g_{i,j}^{(r)}(S)
 :=\Pp(B_{i,j}\mid S^{(r)}=S),
 \qquad S\subseteq\cB_{i,j},
\]
so $g_{i,j}=g_{i,j}^{(x)}$.  A permutation of the labels $x,y,z$ that
sends $x$ to $r$ preserves both the five-state law and $B_{i,j}$ and sends
$S^{(x)}$ to $S^{(r)}$.  Hence, for every
$S\subseteq\cB_{i,j}$,
\[
 \Pp(B_{i,j}\cap\{S^{(r)}=S\})
 =\Pp(B_{i,j}\cap\{S^{(x)}=S\}).
\]
Every state has marginal law $\mu_{\cB_{i,j}}$, so
\[
 \Pp(S^{(r)}=S)=\Pp(S^{(x)}=S)
 =\mu_{\cB_{i,j}}(S)>0.
\]
Dividing the preceding joint-probability identity by this common marginal
gives
\[
 g_{i,j}^{(r)}(S)=g_{i,j}^{(x)}(S)=g_{i,j}(S).
\]
Bayes' formula therefore gives, for each of the three possible states,
\begin{equation}\label{eq:all-state-density}
 \frac{\dd\,\operatorname{Law}(S^{(r)}\mid B_{i,j})}
      {\dd\mu_{\cB_{i,j}}}(S)
 =\frac{g_{i,j}(S)}{\beta_{i,j}}.
\end{equation}

Relative to the product support law on all $M$ bands in group $i$, the
word-selected state density is consequently exactly
\begin{equation}\label{eq:one-group-density}
 G_i=\frac1M\sum_{j=1}^M\frac{g_{i,j}}{\beta_{i,j}}.
\end{equation}
Indeed, conditional on $J_i=j$, every inactive band has density one, while
the active state has density $g_{i,j}/\beta_{i,j}$ by
\eqref{eq:all-state-density}; averaging over $j$ gives
\eqref{eq:one-group-density}.

Under the ambient product law, the $M$ summands in
\eqref{eq:one-group-density} are independent, have mean one, and have
second moment at most $K$.  Cross terms vanish after centering, so
\begin{equation}\label{eq:average-variance}
 \E(G_i-1)^2
 =\frac1{M^2}
 \sum_{j=1}^M
 \E\left(\frac{g_{i,j}}{\beta_{i,j}}-1\right)^2
 =\frac1{M^2}
 \sum_{j=1}^M
 \left[
  \E\left(\frac{g_{i,j}}{\beta_{i,j}}\right)^2-1
 \right]
 \le\frac KM.
\end{equation}
The groups are independent.  Therefore the density of the full word is
\[
 G_\omega=\prod_{i=1}^H G_i.
\]
The telescoping identity
\[
 \prod_{i=1}^HG_i-1
 =
 \sum_{i=1}^H(G_i-1)\prod_{k<i}G_k
\]
together with independence, $\E G_k=1$, and Cauchy--Schwarz gives
\begin{align*}
 \E|G_\omega-1|
 &\le
 \sum_{i=1}^H
 \E|G_i-1|\prod_{k<i}\E G_k\\
 &\le
 \sum_{i=1}^H\sqrt{\E(G_i-1)^2}
 \le H\sqrt{\frac KM}.
\end{align*}
This proves (ii).

It remains to check balance.  On an active band $(i,j)$ of scale
$T_{i,j}$, after
removing the common $c$-contribution, the three normalized logarithmic
contributions are
\[
 Y+Z,\qquad X+Z,\qquad X+Y.
\]
Their pairwise differences are $X-Y,X-Z,Y-Z$.  On $B_{i,j}$ each has
absolute value at most $w_{i,j}$, so the corresponding actual logarithms
have range at most
$T_{i,j}w_{i,j}=\eta_{i,j}\le\eta_*$.  Two words differ in at most this amount in each of
the $H$ coordinates; hence all word logarithms have range at most
$H\eta_*$.  Their arithmetic mean is the logarithm of the geometric mean
of the word products and lies between the minimum and maximum.  This proves
\eqref{eq:flattened-balance}.
\end{proof}

\section{Band placement and completion of the proof}
\label{sec:completion}

We now place finitely many candidate bands inside the prime set used in the
squarefree capacity.  The order of limits is
\[
 H\ \text{first},\qquad M\ \text{second},\qquad
 y\to\infty\ \text{last}.
\]

For $y>\e^\e$, put
\begin{equation}\label{eq:final-prime-set}
 E_y:=\{p:\log y<p\le y\},
 \qquad L:=\log y.
\end{equation}
Fix $H,M$, and for $1\le j\le HM$ define
\begin{equation}\label{eq:band-exponents}
 b_j:=\frac12\vartheta^{\,2(j-1)}
\end{equation}
and
\begin{equation}\label{eq:placed-bands}
 \cB_j:=
 \left\{
 p:\e^{L^{\vartheta b_j}}<p\le\e^{L^{b_j}}
 \right\}.
\end{equation}
Assign bands $(i-1)M+1,\ldots,iM$ to the $i$th coordinate group.

These bands are pairwise disjoint.  Indeed,
\[
 b_{j+1}=\vartheta^2b_j<\vartheta b_j,
\]
so the upper endpoint of $\cB_{j+1}$ is smaller than the lower endpoint of
$\cB_j$.  They also lie in $E_y$ for all sufficiently large $y$.  Their
upper endpoints are below $y$ because $b_j\le1/2<1$, while for the smallest
fixed exponent $b_{HM}>0$,
\[
 L^{\vartheta b_{HM}}>\log L
\]
eventually.  Thus the lower endpoint of every band exceeds
$\e^{\log L}=L=\log y$.

Use the common balancing parameter
\begin{equation}\label{eq:final-eta}
 \eta_y:=\frac1{\log\log y}.
\end{equation}
For band $j$, its scale in \eqref{eq:band-definitions} is
$T_j=L^{b_j}$; define
\begin{equation}\label{eq:placed-window-width}
 w_j:=\frac{\eta_y}{T_j}.
\end{equation}
Thus, in the notation of Proposition~\ref{prop:alternative-bands},
\begin{equation}\label{eq:placed-group-indexing}
 \begin{aligned}
  \cB_{i,j}&:=\cB_{(i-1)M+j},&
  T_{i,j}&:=T_{(i-1)M+j},\\
  w_{i,j}&:=w_{(i-1)M+j},&
  \eta_{i,j}&:=\eta_y
 \end{aligned}
 \qquad(1\le i\le H,\ 1\le j\le M).
\end{equation}
The slow-decay condition holds because
\[
 \eta_y=\frac{b_j}{\log T_j}\longrightarrow0,
 \qquad
 \frac{\log(1/\eta_y)}{\log T_j}
 =
 \frac{\log\log\log y}{b_j\log\log y}
 \longrightarrow0.
\]
Since $HM$ is fixed and $b_{HM}>0$, one threshold
$y_0(H,M)$ makes every estimate in
Sections~\ref{sec:firstmoment}--\ref{sec:secondmoment} valid simultaneously
on all candidate bands; increase this threshold, if necessary, so that
$H\eta_y\le1$, as required by Proposition~\ref{prop:joint-prefix}.
Proposition~\ref{prop:rooted-balance} supplies the
same constant $K$ for every band.  Moreover, on band $j$,
\[
 T_jw_j=\eta_y,
\]
so its actual logarithmic balance error is at most $\eta_y$.

Apply Proposition~\ref{prop:alternative-bands} on the finite set $E_y$.
Primes in inactive bands and outside all candidate bands are sampled into
the common support with their product probabilities, exactly as in that
proposition.  Thus the ambient word law is $\mu_{E_y}$.  The resulting cube
law has
\[
 \eps_{H,M}=H\sqrt{\frac KM},
 \qquad
 \zeta_y=H\eta_y
\]
in Proposition~\ref{prop:joint-prefix}.  Hence
\begin{equation}\label{eq:final-capacity-bound}
 \frac{I_{E_y}^{\mathrm{sf}}}{Z_{E_y}}
 \le
 3\kappacap^H+
 C\left(
 H\sqrt{\frac KM}+H\eta_y
 \right)
\end{equation}
for an absolute constant $C$.

We now execute the stated order of limits.  Given $\eps>0$, first choose
$H$ so that
\[
 3\kappacap^H<\frac{\eps}{3};
\]
this is possible because $\kappacap<1$.  Keeping $H$ fixed, choose $M$ so
large that
\[
 CH\sqrt{\frac KM}<\frac{\eps}{3}.
\]
Finally keep $H,M$ fixed and let $y\to\infty$ until
$CH\eta_y<\eps/3$.  Equation~\eqref{eq:final-capacity-bound} then gives,
for every sufficiently large $y$,
\[
 0\le\frac{I_{E_y}^{\mathrm{sf}}}{Z_{E_y}}<\eps.
\]
Thus the limsup is at most $\eps$.  Since $\eps>0$ was arbitrary and the
ratio is nonnegative, we have proved
\begin{equation}\label{eq:capacity-vanishes}
 \frac{I_{E_y}^{\mathrm{sf}}}{Z_{E_y}}\longrightarrow0.
\end{equation}

The least prime in $E_y$ exceeds $\log y\to\infty$.  Thus
\eqref{eq:capacity-vanishes} verifies both hypotheses of
Corollary~\ref{cor:capacity-target}.  That corollary yields
$f(N)=o(N)$ and completes the proof of Theorem~\ref{thm:main}.

\section{Conclusion}

The proof separates the problem into three reusable interfaces.  The
finite-prime envelope and unit-exponent deletion convert integer density
into a normalized squarefree capacity.  Joint-prefix transference converts
flat balanced pair-product cubes into a cap-set saving.  Finally, the
prime-band construction provides those cubes: the Poisson anchor argument
gives the first moment, the discrete canonical pivots give the matching
second moment, and alternative bands flatten the conditioned marginals.
The order $H$, then $M$, then $y$ closes these interfaces without requiring
any quantitative rate for $f(N)/N$.

\appendix
\section{Numerical Verification}\label{app:verification}

A companion numerical verifier is supplied with the paper as
\texttt{numerical\_verifier.py}.  It can be run with
\begin{center}
\texttt{python3 numerical\_verifier.py}.
\end{center}

The verifier certifies the explicit finite numerical and combinatorial
comparisons appearing in the proof.  More precisely, it verifies:
\begin{itemize}
 \item the squarefree product-law and unit-deletion normalizations;
 \item the substitution in the cap-set constant;
 \item the explicit structural, survival, buffer, and anchor constants;
 \item the five-state and conditional missing-petal probabilities;
 \item atom-by-atom finite tests of factorial insertion;
 \item all nine-mark rank counts and all noncollinear $2\times2$ mark
       matrices;
 \item the elementary balance identities and separation of the candidate
       bands.
\end{itemize}

The program uses only the Python standard library.  Rational identities are
checked exactly with \texttt{Fraction} arithmetic, and comparisons involving
$\log 3$ use a rigorous rational interval obtained from the
$\operatorname{atanh}$ series with an explicit remainder bound.  No binary
floating-point arithmetic or random simulation is used.  The optional flag
\texttt{--self-test} additionally exhausts small valuation models,
squarefree set systems, and pair-product cubes as regression checks.

\begin{thebibliography}{99}
\raggedright

\bibitem{AbbottGardner}
H.~L. Abbott and B.~Gardner,
\newblock An extremal problem in number theory,
\newblock \emph{Canadian Mathematical Bulletin} \textbf{10} (1967),
173--177.
\newblock \url{https://doi.org/10.4153/CMB-1967-015-8}

\bibitem{Bloom536}
T.~F. Bloom,
\newblock \emph{Erd\H{o}s Problem \#536},
\newblock Erd\H{o}s Problems, accessed 13 July 2026.
\newblock \url{https://www.erdosproblems.com/536}

\bibitem{EllenbergGijswijt}
J.~S. Ellenberg and D.~Gijswijt,
\newblock On large subsets of $\mathbb F_q^n$ with no three-term arithmetic
progression,
\newblock \emph{Annals of Mathematics} \textbf{185} (2017), 339--343.
\newblock \url{https://doi.org/10.4007/annals.2017.185.1.8}

\bibitem{Erdos1964}
P.~Erd\H{o}s,
\newblock On a problem in elementary number theory and a combinatorial
problem,
\newblock \emph{Mathematics of Computation} \textbf{18} (1964), 644--646.
\newblock \url{https://doi.org/10.2307/2002950}

\bibitem{IwaniecKowalski}
H.~Iwaniec and E.~Kowalski,
\newblock \emph{Analytic Number Theory},
\newblock American Mathematical Society Colloquium Publications, vol.~53,
American Mathematical Society, Providence, RI, 2004.

\bibitem{MontgomeryVaughan}
H.~L. Montgomery and R.~C. Vaughan,
\newblock \emph{Multiplicative Number Theory I: Classical Theory},
\newblock Cambridge Studies in Advanced Mathematics, vol.~97,
Cambridge University Press, Cambridge, 2007.

\bibitem{MontgomeryVaughanLargeSieve}
H.~L. Montgomery and R.~C. Vaughan,
\newblock The large sieve,
\newblock \emph{Mathematika} \textbf{20} (1973), 119--134.
\newblock \url{https://doi.org/10.1112/S0025579300004708}

\end{thebibliography}

\end{document}
